\chapter{Fundamentals}
\section{Basic Properties of Sets} \label{sec:basic-properties-of-sets}

Many properties of sets follow directly from the axioms listed in \cref{ch:zfc-axioms}. Note that the
although the axioms postulate the existence of various sets (e.g: the empty set, the power set and
the union of any given set), they say nothing about the \emph{uniqueness} of those sets.
Therefore, we shall begin this chapter by proving the uniqueness of the various sets postulated by the basic axioms.
Common notations used in set theory will also be introduced along with a rigorous justification for their usage.

\begin{thm}
    \label{thm:uniqueness-empty-set}
    The empty set is unique.
    \begin{equation*}
        \forall A,B \Big[\, \forall x \, \mleft( x \notin A \mright) \land \forall x \, \mleft( x \notin B \mright) \implies A = B \,\Big]
    \end{equation*}
\end{thm}

\begin{proof}
    Let \( A \) and \( B \) be arbitrary sets that do not contain any elements. That is,
    \begin{equation}
        \label{eq:A-B-empty}
        \forall x \left(x \notin A \right) \land \forall x \left(x \notin B \right)
    \end{equation}
    Let \( x_{0} \) be any arbitrary object. Then by \eqref{eq:A-B-empty}, it follows that \( x_{0} \notin A \land x_{0} \notin B \).
    Thus, the statements \( x_{0} \in A \implies x_{0} \in B \) and \( x_{0} \in B \implies x_{0} \in A \) are both vacuously true. By elementary
    logic,
    \begin{equation*}
        \mleft( x_{0} \in A \implies x_{0} \in B \mright) \land
        \mleft( x_{0} \in B \implies x_{0} \in A \mright) \iff
        \mleft( x_{0} \in A \iff x_{0} \in B \mright)
    \end{equation*}
    Since \( x_{0} \) is arbitrary, we conclude by Universal Generalization that,
    \begin{equation*}
        \forall x \, \mleft( x \in A \iff x \in B \mright)
    \end{equation*}
    By the \nameref{axm:axiom-of-extensionality}, it follows that \( A = B \). Thus, the empty set is unique.
\end{proof}

\vspace{0.3cm}
\noindent Since the empty set is unique, we can denote it with a single universal symbol.
\begin{dfn}
    \label{dfn:empty-set}
    \( \emptyset \) is the empty set.
    \begin{equation*}
        \forall S \, \mleft( S = \emptyset \iff \forall x \, \mleft( x \notin S \mright) \mright)
    \end{equation*}
\end{dfn}

\begin{thm}[Restricted Comprehension]
    \label{thm:uniqueness-predicate-set}
    Given any predicate \( \varphi \) and any set \( A \), the set \( S \) such that
    \( \forall x \, \mleft( x \in S \iff x \in A \land \varphi(x) \mright) \) is unique.
\end{thm}

\begin{proof}
    Let \( S \) and \( S^{\prime} \) be sets such that,
    \begin{gather}
        \forall x \, \mleft( x \in S \iff x \in A \land \varphi(x) \mright) \label{eq:S-satisfies-phi} \\
        \forall x \, \mleft( x \in S^{\prime} \iff x \in A \land \varphi(x) \mright) \label{eq:S-prime-satisfies-phi}
    \end{gather}
    By elementary logic, \eqref{eq:S-satisfies-phi} and \eqref{eq:S-prime-satisfies-phi} clearly imply,
    \begin{equation*}
        \forall x \, \mleft( x \in S \iff x \in S^{\prime} \mright)
    \end{equation*}
    By the \nameref{axm:axiom-of-extensionality}, it follows that \( S = S^{\prime} \).
\end{proof}

\begin{thm}
    \label{thm:sufficient-condition-for-unrestricted-comprehension}
    For any predicate \( \varphi \),
    \begin{equation*}
        \exists A \, \forall x \, \big( \varphi(x) \implies x \in A \big)
        \implies
        \exists! S \, \forall x \, \mleft( x \in S \iff \varphi(x) \mright)
    \end{equation*}
\end{thm}

\begin{proof}
    Let \( \varphi \) be a predicate such that,
    \begin{equation}
        \label{eq:phi-x-implies-x-in-A}
        \forall x \, ( \varphi(x) \implies x \in A )
    \end{equation}
    for some set \( A \).
    The \nameref{axm:axiom-schema-of-specification} and \cref{thm:uniqueness-predicate-set}
    guarantee the existence of a set \( S \) such that,
    \begin{equation}
        \label{eq:property-of-S}
        \forall x \, \mleft( x \in S \iff x \in A \land \varphi(x) \mright)
    \end{equation}
    Let \( x_{0} \in S \) be arbitrary. Then \( x_{0} \,{\in}\, A \land \varphi(x_{0}) \) follows immediately
    from \eqref{eq:property-of-S}. By Conjunction Elimination, we can infer \( \varphi(x_{0}) \).
    \begin{equation}
        \label{eq:x-in-S-implies-phi-x}
        \therefore \, x_{0} \in S \implies \varphi(x_{0})
    \end{equation}
    To prove the converse, let \( x_{0} \) be an arbitrary object such that \( \varphi(x_{0}) \) holds.
    By \eqref{eq:phi-x-implies-x-in-A}, it follows that \( x_{0} \in A \).
    And by Conjunction Introduction, we infer \( x_{0} \in A \land \varphi(x_{0}) \),
    which in turns implies \( x_{0} \in S \) by \eqref{eq:property-of-S}.
    \begin{equation}
        \label{eq:phi-x-implies-x-in-S}
        \therefore \, \varphi(x_{0}) \implies x_{0} \in S
    \end{equation}
    By \eqref{eq:phi-x-implies-x-in-S}, \eqref{eq:x-in-S-implies-phi-x}, and Universal Generalization,
    \begin{equation*}
        \forall x \, \mleft( x \in S \iff \varphi(x) \mright)
    \end{equation*}
    Thus, \( \exists  S \, \forall x \, \mleft( x \in S \iff \varphi(x) \mright) \) by Existential Generalization.
    The proof for the uniqueness of \( S \) is straightforward.
    Suppose that \( S \) and \( S^{\prime} \) are sets such that,
    \begin{gather*}
        \forall x \, \mleft( x \in S \iff \varphi(x) \mright) \\
        \forall x \, \mleft( x \in S^{\prime} \iff \varphi(x) \mright)
    \end{gather*}
    It is obvious that \( \forall x \, \mleft( x \in S \iff x \in S^{\prime} \mright) \) and thus \( S = S^{\prime} \)
    by the \nameref{axm:axiom-of-extensionality}.
\end{proof}

\vspace{0.3cm}
\noindent\cref{thm:uniqueness-predicate-set,thm:sufficient-condition-for-unrestricted-comprehension}
serve as the rigorous foundation for the familiar set-builder notation.
\begin{dfn}[Set-Builder Notation]
    \label{dfn:set-builder-notation}
    For any set \( A \) and any predicate \( \varphi \), the set
    of all elements of \( A \) that satisfy \( \varphi \) is denoted by
    \setbuilder{x \in A}{ \varphi(x) }.
    \begin{equation*}
        \forall A \, \forall x \, \Big[ x \in \setbuilder{z \in A}{\varphi(z)} \iff x \in A \land \varphi(x) \Big]
    \end{equation*}
    Similarly, \setbuilder{x}{\varphi(x)} denotes the set of all \( x \) satisfying \( \varphi(x) \).
    \begin{equation*}
        \forall x \, \Big[ x \in \setbuilder{z}{\varphi(z)} \iff \varphi(x) \Big]
    \end{equation*}
\end{dfn}

\begin{rem}
    The set \setbuilder{x}{\varphi(x)} doesn't exist for every predicate \( \varphi \) (see \cref{thm:russell-paradox}).
    Unless one can prove its existence (with \cref{thm:sufficient-condition-for-unrestricted-comprehension}) first, the notation
    \setbuilder{x}{\varphi(x)} is meaningless and therefore should not be used.
\end{rem}

\begin{thm}
    \label{thm:equal-predicate-specification}
    For any set \( A \) and any two predicates \( \varphi_{1} \) and \( \varphi_{2} \),
    \begin{equation*}
        \forall x \, \big[ \varphi_{1}(x) \iff \varphi_{2}(x) \big]
        \implies \setbuilder{x \in A}{\varphi_{1}(x)} = \setbuilder{x \in A}{\varphi_{2}(x)}
    \end{equation*}
\end{thm}

\begin{proof}
    Let \( A \) be any set and \( \varphi_{1} \) and \( \varphi_{2} \) be two predicates. Suppose that,
    \begin{equation}
        \label{eq:phi1-iff-phi2}
        \forall x \, \big[ \varphi_{1}(x) \iff \varphi_{2}(x) \big]
    \end{equation}
    Let \( S_{1} = \setbuilder{x \in A}{\varphi_{1}(x)} \) and \( S_{2} = \setbuilder{x \in A}{\varphi_{2}(x)} \).
    Let \( x_{0} \in S_{1} \) be arbitrary. By \cref{dfn:set-builder-notation}, it follows that \( x_{0} \in A \land \varphi_{1}(x_{0}) \).
    Thus, \( x_{0} \in A \land \varphi_{2}(x_{0}) \) and \( x_{0} \in S_{2} \) by
    \eqref{eq:phi1-iff-phi2} and \cref{dfn:set-builder-notation}. Since \( x_{0} \) is arbitrary, Universal Generalization yields,
    \begin{equation}
        \label{eq:s1-subset-s2}
        \forall x \, \mleft( x \in S_{1} \implies x \in S_{2} \mright)
    \end{equation}
    Letting \( x_{0} \in S_{2} \) be arbitrary, we can use the same argument in reverse and prove that,
    \begin{equation}
        \label{eq:s2-subset-s1}
        \forall x \, \mleft( x \in S_{2} \implies x \in S_{1} \mright)
    \end{equation}
    Together, \cref{eq:s1-subset-s2,eq:s2-subset-s1} imply the biconditional,
    \begin{equation}
        \label{eq:s1-equals-s2}
        \forall x \, \mleft( x \in S_{1} \iff x \in S_{2} \mright)
    \end{equation}
    \begin{equation*}
        \therefore \, \forall x \, \big[\, \varphi_{1}(x) \iff \varphi_{2}(x) \,\big]
        \implies \setbuilder{x \in A}{\varphi_{1}(x)} = \setbuilder{x \in A}{\varphi_{2}(x)} \qedhere
    \end{equation*}
\end{proof}

\vspace{0.5cm}
\begin{rem}
    Note that the converse of \cref{thm:equal-predicate-specification}
    is \keyword{not} true as illustrated in the following counterexample.
    Consider the predicates,
    \begin{gather}
        \varphi_{1}(x) \coloneq
        \Big[\, \big( \Gamma(x) \land x \in A \big)
        \lor x \in B \,\Big] \label{eq:phi1-definition} \\
        \varphi_{2}(x) \coloneq
        \Big[\, \big( \Gamma(x) \land x \in A \big)
        \lor x \in C \,\Big] \label{eq:phi2-definition}
    \end{gather}
    where \( \Gamma \) is another predicate and \( B \) and \( C \) are distinct,
    non-empty sets that have no element in common with \( A \). In other words,
    \begin{gather}
        B \neq C \label{eq:B-C-Not-Equal} \\
        \forall x \, \mleft( x \in A \implies x \notin B \mright) \label{eq:A-B-no-common-elements} \\
        \forall x \, \mleft( x \in A \implies x \notin C \mright) \label{eq:A-C-no-common-elements}
    \end{gather}
    We will begin by showing that \( \setbuilder{z \in A}{\varphi_{1}(z)} = \setbuilder{z \in A}{\varphi_{2}(z)} \)
    is true for the predicates \( \varphi_{1} \) and \( \varphi_{2} \) as defined in
    \eqref{eq:phi1-definition} and \eqref{eq:phi2-definition}.
    Following that, we will prove that a biconditional relationship does \emph{not} hold between \( \varphi_{1} \) and
    \( \varphi_{2} \) to establish the counterexample.

    \vspace{0.3cm}
    \noindent
    Let \( x_{0} \) be any element of \setbuilder{z \in A}{\varphi_{1}(z)}.
    Thus, we can make the following chain of inferences:
    \begin{equation*}
        \begin{aligned}
            x_{0} \in \setbuilder{z \in A}{\varphi_{1}(z)}& \\
            x_{0} \in A \land \varphi_{1}(x_{0})& \\
            x_{0} \in A& \\
            x_{0} \notin B& \\
            \varphi_{1}(x_{0})& \\
            \big[\, \Gamma(x_{0}) \land x_{0} \in A \,\big] \lor x_{0} \in B& \\
            \Gamma(x_{0}) \land x_{0} \in A& \\
            \therefore \, \Gamma(x_{0})&
        \end{aligned}
        \text{  }
        \begin{aligned}
            &\quad\text{Assumption}\\
            &\quad\text{\cref{dfn:set-builder-notation}}\\
            &\quad\text{Conjuction Elimination}\\
            &\quad\text{by \eqref{eq:A-B-no-common-elements}}\\
            &\quad\text{Conjuction Elimination}\\
            &\quad\text{by \eqref{eq:phi1-definition}}\\
            &\quad\text{Disjunctive Syllogism}\\
            &\quad\text{Conjunction Elimination}
        \end{aligned}
    \end{equation*}
    Since \( \Gamma(x_{0}) \) holds and \( x_{0} \in A \) by assumption, it follows from
    \eqref{eq:phi2-definition} that \( \varphi_{2}(x_{0}) \) also holds. Thus, \( x_{0} \in
    \setbuilder{z \in A}{\varphi_{2}(z)} \). Since \( x_{0} \) is arbitrary, we infer,
    \begin{equation*}
        \forall x \, \Big[ x \in \setbuilder{z \in A}{\varphi_{1}(z)}
        \implies x \in \setbuilder{z \in A}{\varphi_{2}(z)} \Big]
    \end{equation*}
    The exact same argument can be done in reverse to prove the converse, thereby
    establishing the biconditional,
    \begin{equation*}
        \forall x \, \Big[ x \in \setbuilder{z \in A}{\varphi_{1}(z)}
        \iff x \in \setbuilder{z \in A}{\varphi_{2}(z)} \Big]
    \end{equation*}
    Thus, the definitions of \( \varphi_{1} \) and \( \varphi_{2} \) given by
    \eqref{eq:phi1-definition} and \eqref{eq:phi2-definition} do indeed imply
    \( \{ x \in A \mid \varphi_{1}(x) \} = \{ x \in A \mid \varphi_{2}(x) \} \).
    Now let \( b \) be an arbitrary object such that \( b \in B \land b \notin C \).
    Note that the existence of \( b \) is guaranteed because \( B \) and \( C \) are
    non-empty and \( B \neq C \) by \eqref{eq:B-C-Not-Equal}. We now invoke another
    chain of inferences:
    \begin{equation*}
        \begin{aligned}
            b \in B \land b \notin C& \\
            \varphi_{1}(b)& \\
            b \in A \implies b \notin B& \\
            b \notin A& \\
            \lnot \, \Gamma(b) \lor b \notin A& \\
            \lnot \, ( \Gamma(b) \land b \in A )& \\
            \lnot \, ( \Gamma(b) \land b \in A ) \land b \notin C& \\
            \lnot \, ( ( \Gamma(b) \land b \in A ) \lor b \in C )& \\
            \lnot \, \varphi_{2}(b)& \\
            \varphi_{1}(b) \land \lnot \, \varphi_{2}(b)& \\
            \lnot \, ( \varphi_{1}(b) \implies \varphi_{2}(b) )& \\
            \lnot \, ( \varphi_{1}(b) \implies \varphi_{2}(b) )
            \lor \lnot \, ( \varphi_{2}(b) \implies \varphi_{1}(b) )& \\
            \lnot \, ( \varphi_{1}(b) \implies \varphi_{2}(b)
            \land \varphi_{2}(b) \implies \varphi_{1}(b) )& \\
            \lnot \, ( \varphi_{1}(b) \iff \varphi_{2}(b) )& \\
            \exists \, x [ \lnot \, ( \varphi_{1}(x) \iff \varphi_{2}(x) ) ]& \\
            \therefore \, \lnot \, \forall x \, ( \varphi_{1}(x) \iff \varphi_{2}(x) )&
        \end{aligned}
        \text{  }
        \begin{aligned}
            &\quad\text{Assumption} \\
            &\quad\text{by \eqref{eq:phi1-definition}} \\
            &\quad\text{by \eqref{eq:A-B-no-common-elements} and Universal Instantiation} \\
            &\quad\text{Modus Tollens} \\
            &\quad\text{Disjunction Introduction} \\
            &\quad\text{De Morgan's Law} \\
            &\quad\text{Conjunction Introduction} \\
            &\quad\text{De Morgan's Law} \\
            &\quad\text{by \eqref{eq:phi2-definition}} \\
            &\quad\text{Conjunction Introduction} \\
            &\quad\text{Definition of Conditional} \\
            &\quad\text{Disjunction Introduction} \\
            &\quad\text{De Morgan's Law} \\
            &\quad\text{Definition of Biconditionals} \\
            &\quad\text{Existential Generalization} \\
            &\quad\text{Logical Equivalence}
        \end{aligned}
    \end{equation*}
    Thus, we have arrived at a counterexample for the converse of \cref{thm:equal-predicate-specification}.
    For the predicates \( \varphi_{1} \) and \( \varphi_{2} \) as defined in \eqref{eq:phi1-definition} and
    \eqref{eq:phi2-definition}, \setbuilder{x \in A}{\varphi_{1}(x)} = \setbuilder{x \in A}{\varphi_{2}(x)}
    does \keyword{not} imply \( \forall x \, \big( \varphi_{1}(x) \iff \varphi_{2}(x) \big) \).
\end{rem}

\begin{thm}
    \label{thm:empty-specification}
    For any given predicate \( \varphi \), \( \setbuilder{x \in \emptyset}{\varphi(x)} = \emptyset \).
\end{thm}

\begin{proof}
    Let \( \varphi \) be any predicate. Since \( \emptyset \) is an existing
    set (the \nameref{axm:axiom-of-existence}), it follows from the
    \nameref{axm:axiom-schema-of-specification} that \setbuilder{x \in \emptyset}{\varphi(x)}
    exists. Let \( S = \setbuilder{x \in \emptyset}{\varphi(x)} \).
    By \cref{dfn:set-builder-notation},
    \begin{equation}
        \label{eq:definition-of-S}
        \forall x \, \mleft( x \in S \iff x \in \emptyset \land \varphi(x) \mright)
    \end{equation}
    Let \( x_{0} \) be an arbitrary object. By \cref{dfn:set-builder-notation}, \( x_{0} \notin \emptyset \).
    Thus, \( \lnot \, \mleft(x_{0} \in \emptyset \land \varphi(x_{0})\mright) \) and \( x_{0} \notin S \) by \eqref{eq:definition-of-S}.
    Since \( x_{0} \) is arbitrary, we conclude from Universal Generalization that,
    \begin{equation*}
        \forall x \, \mleft( x \notin S \mright)
    \end{equation*}
    By \cref{dfn:empty-set},
    \begin{equation*}
        S = \setbuilder{x \in \emptyset}{\varphi(x)} = \emptyset \qedhere
    \end{equation*}
\end{proof}

\begin{thm}
    \label{thm:uniqueness-unordered-pair}
    For any two objects \( A \) and \( B \), there is only one set whose elements are exactly \( A \) and \( B \).
\end{thm}

\begin{proof}
    Let \( A \) and \( B \) be any two objects. Let \( S \) and \( S^{\prime} \) be sets whose elements are exactly
    \( A \) and \( B \). In other words,
    \begin{gather*}
        \forall x \, \mleft( x \in S \iff x = A \lor x = B \mright) \\
        \forall x \, \mleft( x \in S^{\prime} \iff x = A \lor x = B \mright)
    \end{gather*}
    By elementary logic,
    \begin{equation*}
        \forall x \, ( x \in S \iff x \in S^{\prime} )
    \end{equation*}
    By the \nameref{axm:axiom-of-extensionality},
    \begin{equation*}
        S = S^{\prime} \qedhere
    \end{equation*}
\end{proof}

\vspace{0.3cm}
\noindent \cref{thm:uniqueness-unordered-pair} yields the following definition.
\begin{dfn}[Unordered-Pair]
    \label{dfn:unordered-pair}
    Given any two objects \( A \) and \( B \), the unique set whose elements are exactly
    \( A \) and \( B \) is denoted by \( \mleft\{ A, B \mright\} \).
    \begin{equation*}
        \forall A,B \, \forall x \, \Big[ x \in \mleft\{ A, B \mright\} \iff x = A \lor x = B \Big]
    \end{equation*}
    The set \( \mleft\{ A, B \mright\} \) is called an \keyword{unordered-pair}. If \( A = B \),
    then the notation can be simplied to \( \mleft\{ A, A \mright\} = \mleft\{ A \mright\} \).
    The set \( \mleft\{ A \mright\} \) is known as a \keyword{singleton}\footnote{A set with only a single element}
    (or a \keyword{unit set}).
\end{dfn}

\begin{rem}
    In general, a comma-separated list between curly braces denotes a set whose elements
    are exactly those objects specified in the list. For instance, \( \mleft\{ A,B,C,D,E \mright\} \) denotes
    a set whose elements are exactly \( A \), \( B \), \( C \), \( D \), and \( E \). The \nameref{axm:axiom-of-pairing}
    can be used to justify the existence of any set whose elements can be enumerated this way.
\end{rem}

\begin{thm}
    \label{thm:no-set-is-an-element-of-itself}
    For any set \( A \), \( A \notin A \).
    \begin{equation*}
        \forall A \, \big( A \notin A \big)
    \end{equation*}
\end{thm}

\begin{proof}
    Let \( A \) be an arbitrary set. By the \nameref{axm:axiom-of-pairing}, the singleton \( \mleft\{ A,A \mright\} = \mleft\{ A \mright\} \)
    exists. Clearly, \( A \in \mleft\{ A \mright\} \), so \( \exists x \, \mleft( x \in \mleft\{ A \mright\} \mright) \).
    Thus, by the \nameref{axm:axiom-of-regularity},
    \begin{equation}
        \label{eq:consequence-of-regularity}
        \exists x \, \Big[ x \in \mleft\{ A \mright\} \land \forall y \, \big( y \in x \implies y \notin \mleft\{ A \mright\} \big) \Big]
    \end{equation}
    But by \cref{dfn:unordered-pair}, \( \forall x \, \big( x \in \mleft\{ A \mright\} \iff x = A \big) \), so \eqref{eq:consequence-of-regularity}
    becomes,
    \begin{equation*}
        \exists x \, \Big[ x = A \land \forall y \, \big( y \in x \implies y \neq A \big) \Big]
    \end{equation*}
    By \cref{logic:using-existential-quantifier-to-express-the-fact-that-an-object-satisfies-a-predicate}, this is equivalent to,
    \begin{equation*}
        \forall y \, \big( y \in A \implies y \neq A \big) \iff \forall y \, \big( y = A \implies y \notin A \big)
    \end{equation*}
    By Universal Instantiation, we thus have,
    \begin{equation*}
        A = A \implies A \notin A
    \end{equation*}
    Since \( A = A \) is trivially true, we conclude \( A \notin A \) by Modus Ponens. Since \( A \) is arbitrary, the theorem follows immediately
    from Universal Generalization.
\end{proof}

\begin{thm}
    \label{thm:uniqueness-power-set}
    For any set \( A \), there is only one set \( P \) such that,
    \begin{equation*}
        \forall x \mleft( x \in P
        \iff \forall y \mleft( y \in x \implies y \in A \mright) \mright)
    \end{equation*}
\end{thm}

\begin{proof}
    Let \( A \) be an arbitrary set. Let \( P \) and \( P^{\prime} \) be two sets
    such that,
    \begin{gather}
        \forall x \, ( x \in P \iff \forall y \, ( y \in x \implies y \in A ) ) \label{eq:p-biconditional} \\
        \forall x \, ( x \in P^{\prime} \iff \forall y \, ( y \in x \implies y \in A ) ) \label{eq:p-prime-biconditional}
    \end{gather}
    By elementary logic, \eqref{eq:p-biconditional} and \eqref{eq:p-prime-biconditional} implies,
    \begin{gather*}
        \forall x \, ( x \in P \iff x \in P^{\prime} ) \\
        \therefore \, P = P^{\prime} \qedhere
    \end{gather*}
\end{proof}

\vspace{0.3cm}
\noindent\cref{thm:uniqueness-power-set} justifies the following definition of power sets.

\begin{dfn}[Power Set]
    \label{dfn:power-set}
    For any set \( A \), the \keyword{power set} of \( A \) is given by,
    \begin{equation*}
        \powerset{A} = \setbuilder{S}{\forall x \, \mleft( x \in S \implies x \in A \mright)}
    \end{equation*}
\end{dfn}

\begin{thm}
    \label{thm:power-set-is-always-non-empty}
    For any set \( A \), \( \powerset{A} \neq \emptyset \).
    \begin{equation*}
        \forall A \, ( \powerset{A} \neq \emptyset )
    \end{equation*}
\end{thm}

\begin{proof}
    Let \( A \) be an arbitrary set. By \cref{thm:every-set-is-subset-of-itself}, it follows that
    \( A \subseteq A \). Thus,
    \begin{equation*}
        \begin{aligned}
            A \in \powerset{A}& \\
            \exists x ( x \in \powerset{A} )& \\
            \lnot \, \forall x \, ( x \notin \powerset{A} )&\\
            \powerset{A} \neq \emptyset&
        \end{aligned}
        \text{  }
        \begin{aligned}
            &\quad\text{\cref{dfn:power-set}}\\
            &\quad\text{Existential Generalization}\\
            &\quad\text{Logical Equivalence}\\
            &\quad\text{\cref{dfn:empty-set}}\\
        \end{aligned}
    \end{equation*}
    Since \( A \) is arbitrary, the theorem follows from Universal Generalization.
\end{proof}

\begin{thm}
    \label{thm:empty-set-belongs-to-the-powerset-of-any-set}
    For any set \( A \), \( \emptyset \in \powerset{A} \).
    \begin{equation*}
        \forall A \, \mleft( \emptyset \in \powerset{A} \mright)
    \end{equation*}
\end{thm}

\begin{proof}
    Let \( A \) be an arbitrary set. By \cref{dfn:power-set},
    \begin{equation}
        \label{eq:apply-definition-of-power-set}
        \forall x \, \mleft( x \in \powerset{A} \iff
        \forall y \, \mleft( y \in x \implies y \in A \mright) \mright)
    \end{equation}
    But \( \forall x \, \mleft( x \notin \emptyset \mright) \) by
    \cref{dfn:empty-set}. Thus, \( y_{0} \in \emptyset \implies y_{0} \in A \)
    is vacuously true for any arbitrary \( y_{0} \).
    By Universal Generalization,
    \begin{equation*}
        \forall y \, \mleft( y \in \emptyset \implies y \in A \mright)
    \end{equation*}
    By \eqref{eq:apply-definition-of-power-set} and Universal Instantiation,
    \begin{equation*}
        \emptyset \in \powerset{A}
    \end{equation*}
    Since \( A \) is arbitrary, the theorem follows from Universal Generalization.
\end{proof}

\begin{thm}
    \label{thm:every-set-is-an-element-of-its-powerset}
    For any set \( A \), \( A \in \powerset{A} \)
    \begin{equation*}
        \forall A \, \mleft( A \in \powerset{A} \mright)
    \end{equation*}
\end{thm}

\begin{proof}
    See \cref{thm:power-set-is-always-non-empty}.
\end{proof}

\begin{thm}[Russell's Paradox]
    \label{thm:russell-paradox}
    The set \setbuilder{x}{x \notin x} does not exist.
\end{thm}

\begin{proof}
    Suppose that \( R = \setbuilder{x}{x \notin x} \) exists. Thus,
    by \cref{dfn:set-builder-notation},
    \begin{equation}
        \label{eq:property-of-russell-set}
        \forall x \, \mleft( x \in R \iff x \notin x \mright)
    \end{equation}
    But \eqref{eq:property-of-russell-set} implies that
    \( R \in R \iff R \notin R \), which is a logical
    impossiblity. The supposition that \setbuilder{x}{x \notin x}
    exists must therefore be false.
\end{proof}

\begin{thm}
    \label{thm:the-set-of-all-sets-does-not-exist}
    The set of all sets does not exist.
    \begin{equation*}
        \forall A \, \exists x \, \mleft( x \notin A \mright)
    \end{equation*}
\end{thm}

\begin{proof}
    Suppose that the set of all sets exists. That is,
    \begin{equation}
        \label{eq:set-of-all-sets-property}
        \forall x \, \mleft( x \in A \mright)
    \end{equation}
    for some set \( A \). Let the predicate \( \varphi \)
    be defined as:
    \begin{equation}
        \label{eq:predicate-for-russell-set}
        \varphi(x) \coloneq x \notin x
    \end{equation}
    By the \nameref{axm:axiom-schema-of-specification}, the set
    \( S = \setbuilder{x \in A}{\varphi(x)} = \setbuilder{x \in A}{x \notin x} \)
    exists.
    And by \cref{dfn:set-builder-notation},
    \begin{equation}
        \label{eq:S-in-predicate-form}
        \forall x \, \mleft( x \in S \iff x \in A \land x \notin x \mright)
    \end{equation}
    From elementary logic\footnote{The Law of Excluded Middle}, we know that,
    \begin{equation}
        \label{eq:tautology-rule-S-in-S-or-S-notin-S}
        S \in S \lor S \notin S
    \end{equation}
    If \( S \notin S \), then \( \varphi(S) \) must hold by
    \eqref{eq:predicate-for-russell-set}. And by \eqref{eq:set-of-all-sets-property},
    it follows that \( S \in A \). Thus, \( S \,{\in}\, A \,\,{\land}\,\, \varphi(S) \), so
    \( S \,{\in}\, S \) by \eqref{eq:S-in-predicate-form}. We thus have \( S \,{\notin}\, S \,{\implies}\,
    S \,{\in}\, S \), a logical impossibility. By Reduction ad Absurdum, we conclude that
    \( S \notin S \) cannot possibly be true. Thus, by Disjunctive Syllogism, we conclude that,
    \begin{equation*}
        S \in S
    \end{equation*}
    But if \( S \in S \), then \( S \notin S \) by \eqref{eq:S-in-predicate-form}, yet another
    contradiction. In other words, if the set of all sets exists, then \eqref{eq:tautology-rule-S-in-S-or-S-notin-S}
    had to be false, which is impossible by the Law of Excluded Middle.
    Thus, the set of all sets does not exist.
\end{proof}

\section{Elementary Set Relations and Operations}
The basic properties established in \cref{sec:basic-properties-of-sets} can be used
to rigorously define the familiar set relations and operations.
The process involved in these definitions is straightforward. We first choose a
predicate, and then prove (using the properties established in
\cref{sec:basic-properties-of-sets}) that it determines a unique set. Finally, we use
the predicate to define the set operation.

\subsection{Subset Relation}
\begin{dfn}[Subset]
    \label{dfn:subset}
    For any given set \( A \), a set \( B \) is a \keyword{subset} of \( A \) \annotation{(denoted as \( B \subseteq A \))}
    if and only if every element of \( B \) is also an element of \( A \).
    \begin{equation*}
        \forall A,B \, \big( B \subseteq A \iff \forall x \, \mleft( x \in B \implies x \in A \mright) \big)
    \end{equation*}
\end{dfn}

\begin{dfn}
    \label{dfn:proper-subset}
    Given any set \( A \), a set \( B \) is a \keyword{proper subset} of A
    \annotation{(denoted as \( B \subset A \))} if and only if \( B \subseteq A \land B \neq A \).
    \begin{equation*}
        \forall A,B \, \mleft( B \subset A \iff B \subseteq A \land B \neq A \mright)
    \end{equation*}
\end{dfn}

\subsection{Union}
\begin{thm}
    \label{thm:uniqueness-union}
    Given any set \( A \), the set \( U \)such that \( \forall x \, \big( x \in U \iff \exists S \, {\in} \, A \, \mleft( x \in S \mright) \big) \)
    is unique.
\end{thm}

\begin{proof}
    Let \( A \) be an arbitrary system of sets. The \nameref{axm:axiom-of-union} postulates the existence of a set \( U \) such that,
    \begin{equation}
        \label{eq:U-defining-statement}
        \forall x \, \big( x \in U \iff \exists S \, {\in} \, A \, \mleft( x \in S \mright) \big)
    \end{equation}
    Suppose that there exists another set \( U^{\prime} \) such that,
    \begin{equation}
        \label{eq:U-prime-defining-statement}
        \forall x \, \big( x \in U^{\prime} \iff \exists S \,{\in}\, A \, \mleft( x \in S \mright) \big)
    \end{equation}
    By elementary logic, \eqref{eq:U-defining-statement} and \eqref{eq:U-prime-defining-statement} implies
    \( \forall x \, \mleft( x \in U \iff x \in U^{\prime} \mright) \).
    \begin{equation*}
        \therefore U = U^{\prime} \qedhere
    \end{equation*}
\end{proof}

\vspace{0.3cm}
\noindent \cref{thm:uniqueness-union} gives rise to the following definition.
\begin{dfn}[Union]
    \label{dfn:union-of-sets}
    For any system of sets \( S \), the union of \( S \) is the unique set given by,
    \begin{equation*}
        \Union S = \setbuilder{x}{\exists \, T \, \mleft( T \in S \land x \in T \,\mright)}
    \end{equation*}
    In particular, if \( S = \mleft\{ A, B \mright\} \), then \( \Union S \) can be denoted as \( A \union B \).
\end{dfn}

\begin{thm}
    \label{thm:union-in-terms-of-logical-or}
    Let \( A \) and \( B \) be arbitrary sets. Then \( A \union B = \setbuilder{x}{x \in A \lor x \in B} \). In other words,
    \begin{equation*}
        \forall A,B \, \forall x \, \mleft( x \in A \union B \iff x \in A \lor x \in B \mright)
    \end{equation*}
\end{thm}

\begin{proof}
    Let \( A \) and \( B \) be arbitrary sets. From \cref{dfn:union-of-sets},
    \begin{equation*}
        A \union B = \Union \mleft\{ A, B \mright\} = \setbuilderbig{x}{\exists \, T \, \mleft(\, T \in \mleft\{ A, B \mright\} \land x \in T \,\mright) }
    \end{equation*}
    By \cref{dfn:unordered-pair},
    \begin{equation*}
        T \in \mleft\{ A, B \mright\} \iff \mleft( T = A \mright) \lor \mleft( T = B \mright)
    \end{equation*}
    By elementary logic,
    \begin{align}
        \label{eq:distributive-law-of-logical-and-over-or}
        T \in \mleft\{ A,B \mright\} \land x \in T
        &\iff
        \big( T = A \lor T = B \big) \land x \in T \notag \\
        &\iff
        \big(\, T = A \land x \in T \,\big) \lor \big(\, T = B \land x \in T \,\big)
    \end{align}
    By Existential Generalization,
    \begin{align}
        \label{eq:equivalent-predicate-for-union-of-pair}
        \exists \, T \, \big( T \in \{ A,B \} \land x \in T \big)
        &\iff
        \exists \, T \, \big( T = A \land x \in T \big) \lor
        \exists \, T \, \big( T = B \land x \in T \big) \notag \\
        &\iff x \in A \lor x \in B
    \end{align}
    \eqref{eq:equivalent-predicate-for-union-of-pair} follows readily from the obvious fact that
    \( \exists \, T \, \mleft( T = A \land x \in T \mright) \iff x \in A \) and
    \( \exists \, T \, \mleft( T = B \land x \in T \mright) \iff x \in B \).
    \begin{equation*}
        \therefore \, \forall x \, \Big[\, \exists \, T \, \big( T \in \{ A,B \} \land x \in T \big)
        \iff x \in A \lor x \in B \,\Big]
    \end{equation*}
    Treating each side of the biconditional as a predicate, \cref{thm:equal-predicate-specification}
    yields,
    \begin{equation*}
        A \union B = \setbuilder{x}{\exists \, T \, \big( T \in \left\{ A, B \right\} \land x \in T \big)}
                 = \setbuilder{x}{x \in A \lor x \in B} \qedhere
    \end{equation*}
\end{proof}

\subsection{Intersection}
\begin{thm}
    \label{thm:existence-and-uniqueness-of-intersection}
    For any non-empty system of sets \( A \), there exists a unique set S such that
    \( \forall x \, \mleft( x \in S \iff
    \forall z \, \mleft( z \in A \implies x \in z \mright) \mright) \).
    \begin{equation*}
        \forall A \, \mleft( A \neq \emptyset \implies
        \exists! S \, \forall x \, \mleft( x \in S
        \iff \forall z \, \mleft( z \in A \implies x \in z \mright) \mright) \mright)
    \end{equation*}
\end{thm}

\begin{proof}
    Let \( A \) be an arbitrary, non-empty system of sets.
    \begin{equation}
        \label{eq:A-has-at-least-one-elem}
        \therefore \, \exists x \, \mleft( x \in A \mright)
    \end{equation}
    Let \( \varphi \) be the predicate defined as,
    \begin{equation}
        \label{eq:intersection-predicate-definition}
        \varphi(x) \coloneq \forall z \, \mleft( z \in A \implies x \in z \mright)
    \end{equation}
    We now prove that \( \forall x \, \mleft( \varphi(x) \implies x \in \Union A \mright) \)
    via the following chain of inferences:
    \begin{equation*}
        \begin{aligned}
            \varphi(t)&\\
            \forall z \, \mleft( z \in A \implies t \in z \mright)&\\
            \alpha \in A& \\
            \alpha \in A \implies t \in \alpha & \\
            t \in \alpha& \\
            \alpha \in A \land t \in \alpha& \\
            \exists S \, \mleft( S \in A \land t \in S \mright)& \\
            t \in \union A& \\
            \varphi(t) \implies t \in \union A& \\
            \therefore \, \forall x \, \mleft( \varphi(x) \implies x \in \union A \mright)& \\
        \end{aligned}
        \text{  }
        \begin{aligned}
            &\quad\text{Assumption} \\
            &\quad\text{by \eqref{eq:intersection-predicate-definition}} \\
            &\quad\text{Existential Instantiation, \eqref{eq:A-has-at-least-one-elem}} \\
            &\quad\text{Universal Instantiation} \\
            &\quad\text{Modus Ponens} \\
            &\quad\text{Conjunction Introduction} \\
            &\quad\text{Existential Generalization} \\
            &\quad\text{\cref{dfn:union-of-sets}} \\
            &\quad\text{Deduction Theorem} \\
            &\quad\text{Universal Generalization} \\
        \end{aligned}
    \end{equation*}
    \cref{thm:sufficient-condition-for-unrestricted-comprehension} guarantees the existence
    of the unique set \( S = \setbuilder{x}{\forall z \, \mleft( z \in A \implies x \in z \mright)} \).
    By the Deduction Theorem,
    \begin{equation*}
        A \neq \emptyset \implies \exists! S \, \forall x \, \mleft( x \in S \iff
        \forall z \, \mleft( z \in A \implies x \in z \mright)\mright)
    \end{equation*}
    Since \( A \) is arbitrary, the theorem follows from Universal Generalization.
\end{proof}

\begin{dfn}[Intersection]
    \label{dfn:intersection-of-sets}
    For any non-empty system of sets \( S \), the unique set \( \Intersect S =
    \setbuilder{x}{\forall z \, \mleft( z \in S \implies x \in z \mright)}\) is called
    the \keyword{intersection} of \( S \). If \( S = \mleft\{A, B\mright\} \), we can denote
    \( \Intersect S \) as \( A \intersect B \).
\end{dfn}

\begin{thm}
    \label{thm:intersection-in-terms-of-logical-and}
    For any two sets \( A \) and \( B \), \( A \intersect B = \setbuilder{x}{x \in A \land x \in B} \).
    \begin{equation*}
        \forall A,B \, \forall x \, \mleft( x \in A \intersect B \iff x \in A \land x \in B \mright)
    \end{equation*}
\end{thm}

\begin{proof}
    Let \( A \) and \( B \) be arbitrary sets. By \cref{dfn:intersection-of-sets},
    \begin{equation}
        \label{eq:express-A-intersect-B-using-setbuilder-notation}
        A \intersect B = \Intersect \mleft\{ A, B \mright\} =
        \setbuilder{x}{\varphi(x)}
    \end{equation}
    where,
    \begin{equation}
        \label{eq:introduce-the-predicate-symbol}
        \varphi(x) \coloneq \forall z \, \mleft( z \in
    \mleft\{ A,B \mright\} \implies x \in z \mright)
    \end{equation}
    By \cref{dfn:unordered-pair},
    \begin{equation}
        \label{eq:use-the-definition-of-unordered-pair}
        \forall z \, \mleft( z \in \mleft\{ A,B \mright\} \iff
        z = A \lor z = B \mright)
    \end{equation}
    Let \( x_{0} \) be an arbitary object. Using elementary logic, we reason as follows:
    \begin{equation*}
        \begin{aligned}
            \varphi(x_{0})& \\
            \forall z \, \mleft( z \in \mleft\{ A,B \mright\} \implies x_{0} \in z \mright)& \\
            A \in \mleft\{ A,B \mright\} \implies x_{0} \in A& \\
            B \in \mleft\{ A,B \mright\} \implies x_{0} \in B& \\
            A \in \mleft\{ A,B \mright\} \land B \in \mleft\{ A,B \mright\}& \\
            x_{0} \in A \land x_{0} \in B& \\
            \therefore \, \varphi(x_{0}) \implies \mleft( x_{0} \in A \land x_{0} \in B \mright)& \\
        \end{aligned}
        \text{  }
        \begin{aligned}
            &\quad\text{Assumption}\\
            &\quad\text{by \eqref{eq:introduce-the-predicate-symbol}}\\
            &\quad\text{Universal Instantiation}\\
            &\quad\text{Universal Instantiation}\\
            &\quad\text{by \eqref{eq:use-the-definition-of-unordered-pair}}\\
            &\quad\text{Modus Ponens}\\
            &\quad\text{Deduction Theorem}\\
        \end{aligned}
    \end{equation*}
    Similarly, we can prove the converse via the following chain of reasoning:
    \begin{equation*}
        \begin{aligned}
            x_{0} \in A \land x_{0} \in B& \\
            z_{0} \in \left\{ A,B \right\}& \\
            z_{0} = A \lor z_{0} = B& \\
            z_{0} = A& \\
            x_{0} \in z_{0}& \\
            z_{0} = A \implies x_{0} \in z_{0}& \\
            z_{0} = B& \\
            x_{0} \in z_{0}& \\
            z_{0} = B \implies x_{0} \in z_{0}& \\
            x_{0} \in z_{0}& \\
            z_{0} \in \mleft\{ A,B \mright\} \implies x_{0} \in z_{0}& \\
            \forall z \, \mleft( z \in \mleft\{ A,B \mright\} \implies x_{0} \in z \mright)& \\
            \varphi(x_{0})& \\
            \therefore \, x_{0} \in A \land x_{0} \in B \implies \varphi(x_{0})&
        \end{aligned}
        \text{  }
        \begin{aligned}
            &\quad\text{Assumption}\\
            &\quad\text{Assumption}\\
            &\quad\text{by \eqref{eq:use-the-definition-of-unordered-pair}}\\
            &\quad\text{Assumption}\\
            &\quad\text{\nameref{axm:axiom-of-extensionality}}\\
            &\quad\text{Deduction Theorem}\\
            &\quad\text{Assumption}\\
            &\quad\text{\nameref{axm:axiom-of-extensionality}}\\
            &\quad\text{Deduction Theorem}\\
            &\quad\text{Disjunction Elimination}\\
            &\quad\text{Deduction Theorem}\\
            &\quad\text{Universal Generalization}\\
            &\quad\text{by \eqref{eq:introduce-the-predicate-symbol}}\\
            &\quad\text{Deduction Theorem}
        \end{aligned}
    \end{equation*}
    Combining the results from both analyses via Biconditional Introduction yields,
    \begin{equation*}
        \varphi(x_{0}) \iff x_{0} \in A \land x_{0} \in B
    \end{equation*}
    Since \( x_{0} \) is arbitrary, we conclude by Universal Generalization that,
    \begin{equation*}
        \forall x \, \mleft( \varphi(x) \iff x \in A \land x \in B \mright)
    \end{equation*}
    And since, \( \forall x \, \mleft( x \in A \intersect B \iff \varphi(x) \mright) \), we thus have,
    \begin{equation*}
        A \intersect B = \setbuilder{x}{\varphi(x)} = \setbuilder{x}{x \in A \land x \in B} \qedhere
    \end{equation*}
\end{proof}

\begin{dfn}[Disjoint]
    \label{dfn:disjoint-two-sets}
    Two sets \( A \) and \( B \) are \keyword{disjoint} if and only if \( A \intersect B = \emptyset \).
\end{dfn}

\begin{dfn}[Pairwise Disjoint]
    \label{dfn:pairwise-disjoint-sets}
    A non-empty system of sets \( S \) is \keyword{pairwise disjoint\,\footnote{Some mathematicians use the
    term disjoint as synonymous to pairwise disjoint, but we will avoid that to prevent any ambiguity.
    Instead, we use \keyword{pairwise disjoint} when we're talking about a system of sets and reserve
    the term \keyword{disjoint} for two sets.}} if and only if the intersection of any pair of distinct
    sets from \( S \) is empty. In other words, \( S \) is pairwise disjoint if and only if,
    \begin{equation*}
        \forall A \,{\in}\, S, B \,{\in}\, S \, \mleft( A \neq B \implies A \intersect B = \emptyset \mright)
    \end{equation*}
\end{dfn}

\subsection{Complementation}
\begin{thm}
    \label{thm:existence-and-uniqueness-of-difference}
    For any two sets \( A \) and \( B \), there exists a unique set \( S \) such that
    \( \forall x \, \mleft( x \in S \iff x \in A \land x \notin B \mright) \).
    \begin{equation*}
        \forall A,B \, \exists! S \, \forall x \, \mleft( x \in S \iff x \in A \land x \notin B \mright)
    \end{equation*}
\end{thm}

\begin{proof}
    Let \( A \) and \( B \) be two arbitrary sets. Consider the predicate,
    \begin{equation}
        \label{eq:introduce-the-predicate-for-diff}
        \varphi(x) \coloneq x \in A \land x \notin B
    \end{equation}
    Clearly, \eqref{eq:introduce-the-predicate-for-diff} implies,
    \begin{equation*}
        \forall x \, \mleft( \varphi(x) \implies x \in A \mright)
    \end{equation*}
    Thus, \cref{thm:sufficient-condition-for-unrestricted-comprehension} guarantess the existence of a unique set
    \( S \) such that,
    \begin{gather*}
        \forall x \, \mleft( x \in S \iff x \in A \land x \notin B \mright) \\
        \therefore \, \exists! S \, \forall x \, \mleft( x \in S \iff x \in A \land x \notin B \mright)
    \end{gather*}
    Since \( A \) and \( B \) are arbitrary, the theorem follows from Universal Generalization.
\end{proof}

\begin{dfn}[Relative Complement]
    \label{dfn:difference}
    For any two sets \( A \) and \( B \), the \keyword{complement\,\footnote{Also known as
    the \keyword{set difference}.}} of \( B \) relative to \( A \) is the set given by,
    \begin{equation*}
        A - B = \setbuilder{x}{x \in A \land x \notin B}
    \end{equation*}
    In other words,
    \begin{equation*}
        \forall A,B \, \mleft( x \in \mleft( A - B \mright) \iff x \in A \land x \notin B \mright)
    \end{equation*}
\end{dfn}

\begin{thm}
    \label{thm:existence-and-uniqueness-of-symmetric-difference}
    For any two sets \( A \) and \( B \), there exists a unique set \( S \) such that
    \( \forall x \, \mleft( x \in S \iff \mleft( x \in A \lor x \in B
    \mright) \land \lnot \, \mleft( x \in A \land x \in B \mright) \mright) \).
    \begin{equation*}
        \forall A,B \, \exists! S \, \forall x \, \big( x \in S \iff
        \mleft( x \in A \lor  x \in B \mright) \land
        \lnot \, \mleft( x \in A \land x \in B \mright) \big)
    \end{equation*}
\end{thm}

\begin{proof}
    Let \( A \) and \( B \) be arbitrary sets. Consider the predicate,
    \begin{equation}
        \label{eq:introduce-predicate-for-symmetric-difference}
        \varphi(x) \coloneq \mleft( x \in A \lor x \in B \mright)
        \land
        \lnot \, \mleft( x \in A \land x \in B \mright)
    \end{equation}
    The following chain of inferences should complete the proof:
    \begin{equation*}
        \begin{aligned}
            \varphi(x_{0})& \\
            x_{0} \in A \lor x_{0} \in B& \\
            x_{0} \in A \union B& \\
            \varphi(x_{0}) \implies x_{0} \in A \union B& \\
            \therefore \, \forall x \, \mleft( \varphi(x_{0}) \implies
            x_{0} \in A \union B \mright)&
        \end{aligned}
        \text{  }
        \begin{aligned}
            &\quad\text{Assumption}\\
            &\quad\text{\eqref{eq:introduce-predicate-for-symmetric-difference},
            Conjunction Elimination}\\
            &\quad\text{\cref{thm:union-in-terms-of-logical-or}}\\
            &\quad\text{Deduction Theorem}\\
            &\quad\text{Universal Generalization}
        \end{aligned}
    \end{equation*}
    Thus, by \cref{thm:sufficient-condition-for-unrestricted-comprehension},
    the set \setbuilder{x}{\varphi(x)} exists and is unique.
\end{proof}

\begin{dfn}[Symmetric Difference]
    \label{dfn:symmetric-difference}
    For any two sets \( A \) and \( B \), the \keyword{symmetric difference}
    of \( A \) and \( B \) is the unique set given by,
    \begin{align*}
        A \symdiff B &= \setbuilder{x}{ \mleft( x \in A \lor x \in B \mright) \land
        \lnot \, \mleft( x \in A \land x \in B \mright) } \\
        &=
        \setbuilder{x}{\mleft( x \in A \land x \notin B \mright) \lor \mleft( x \notin A \land x \in B \mright)}
    \end{align*}
\end{dfn}

\section{Set Identities and Laws}
In this section, we present the laws and identities that the basic set-theoretic operations
obey. These identities serves as a means to simplifying and evaluating set-theoretic expressions
and are applicable to many other areas of mathematics where sets of complex mathematical objects
are being studied.

\subsection{Subset Relation}
\begin{thm}
    \label{thm:empty-set-universal-subset}
    For any set \( A \), \( \emptyset \subseteq A \).
    \begin{equation*}
        \forall A \, ( \emptyset \subseteq A )
    \end{equation*}
\end{thm}

\begin{proof}
    Let \( A \) be an arbitrary set. Let \( x \in A \) be arbitrary. By \cref{dfn:empty-set},
    it follows that \( x \notin \emptyset \). Thus, the \( x \in \emptyset \implies x \in A \)
    is vacuously true. By Universal Generalization,
    \begin{equation*}
        \forall x \, \mleft( x \in \emptyset \implies x \in A \mright)
    \end{equation*}
    By \cref{dfn:subset},
    \begin{equation*}
        \emptyset \subseteq A
    \end{equation*}
    Since \( A \) is arbitrary, the theorem follows from Unviersal Generalization.
\end{proof}

\begin{thm}
    \label{thm:every-set-is-subset-of-itself}
    For any set \( A \), \( A \subseteq A \).
    \begin{equation*}
        \forall A \, \mleft( A \subseteq A \mright)
    \end{equation*}
\end{thm}

\begin{proof}
    Let \( A \) be an arbitrary set. Let \( x \) be an arbitrary object.
    Using elementary propositional logic, we can make the following inferences:
    \begin{equation*}
        \begin{aligned}
            x \notin A \lor x \in A& \\
            \mleft( x \in A \mright) \implies \mleft( x \in A \mright)& \\
            \forall x \, \mleft( x \in A \implies x \in A \mright)&
        \end{aligned}
        \text{  }
        \begin{aligned}
            &\quad\text{Law of Excluded Middle}\\
            &\quad\text{Logical Equivalence}\\
            &\quad\text{Universal Generalization}
        \end{aligned}
    \end{equation*}
    \( A \subseteq A \) folllows readily from \cref{dfn:subset}.
\end{proof}

\begin{thm}[Set Equality]
    \label{thm:set-equality-in-terms-of-subsets}
    Two sets are equal if and only if they are subsets of each other.
    \begin{equation*}
        \forall A,B \, \mleft( A = B \iff A \subseteq B \land B \subseteq A \mright)
    \end{equation*}
\end{thm}

\begin{proof}
    Let \( A \) and \( B \) be two arbitrary sets. Suppose that \( A = B \).
    Then by the \nameref{axm:axiom-of-extensionality},
    \begin{gather}
        \forall x \, \mleft( x \in A \iff x \in B \mright) \label{eq:A-B-equality-biconditional} \\
        \therefore \, \forall x \, \mleft( x \in A \implies x \in B \mright)
        \land
        \forall x \, \mleft( x \in B \implies x \in A \mright) \label{eq:A-B-equality-biconditional-conjunction}
    \end{gather}
    By \cref{dfn:subset},
    \begin{gather}
        A \subseteq B \iff \forall x \, \mleft( x \in A \implies x \in B \mright) \label{eq:A-B-equality-rightward-implication} \\
        B \subseteq A \iff \forall x \, \mleft( x \in B \implies x \in A \mright) \label{eq:A-B-equality-leftward-implication}
    \end{gather}
    Applying elementary propositional logic to \eqref{eq:A-B-equality-biconditional-conjunction},
    \eqref{eq:A-B-equality-rightward-implication} and \eqref{eq:A-B-equality-leftward-implication} yields,
    \begin{gather}
        A \subseteq B \land B \subseteq A \notag \\
        \therefore \, A = B \implies \mleft( A \subseteq B \land B \subseteq A \mright) \label{eq:subset-equality-direct}
    \end{gather}
    To prove the converse, suppose that \( A \subseteq B \land B \subseteq A \). Again,
    by \cref{dfn:subset}, we once again arrive at \eqref{eq:A-B-equality-biconditional}; so \( A = B \).
    \begin{equation}
        \label{eq:subset-equality-converse}
        \therefore \, A \subseteq B \land B \subseteq A \implies A = B
    \end{equation}
    Together, \eqref{eq:subset-equality-direct} and \eqref{eq:subset-equality-converse} yields the biconditional,
    \begin{equation*}
        A = B \iff A \subseteq B \land B \subseteq A
    \end{equation*}
    Since \( A \) and \( B \) are arbitrary, the theorem follows from Universal Generalization.
\end{proof}

\begin{thm}
    \label{thm:power-set-of-A-can-never-be-subset-of-A}
    For any set \( A \), \( \powerset{A} \nsubseteq A \).
\end{thm}

\begin{proof}
    Let \( A \) be an arbitrary set. Consider once again the predicate
    defined in \eqref{eq:predicate-for-russell-set} and the set,
    \begin{equation}
        \label{eq:define-S-using-russell-predicate}
        S = \setbuilder{x \in A}{x \notin x}
    \end{equation}
    The following chain of inferences completes the proof:
    \begin{equation*}
        \begin{aligned}
            S \subseteq A& \\
            S \in \powerset{A}& \\
            S \in S& \\
            S \notin S \land S \in A& \\
            S \notin S& \\
            \lnot \, \mleft(S \in A \land S \notin S\mright)& \\
            S \notin A \lor S \in S& \\
            S \notin A& \\
            S \in \powerset{A} \land S \notin A& \\
            \lnot \, \mleft( S \in \powerset{A} \implies S \in A \mright)& \\
            \exists x \, \lnot \, \mleft( x \in \powerset{A} \implies x \in A \mright)& \\
            \lnot \, \forall x \, \mleft( x \in \powerset{A} \implies x \in A \mright)&
        \end{aligned}
        \text{  }
        \begin{aligned}
            &\quad\text{by \eqref{eq:define-S-using-russell-predicate}} \\
            &\quad\text{\cref{dfn:power-set}} \\
            &\quad\text{Assumption} \\
            &\quad\text{by \eqref{eq:define-S-using-russell-predicate}}\\
            &\quad\text{Reduction ad Absurdum}\\
            &\quad\text{by \eqref{eq:define-S-using-russell-predicate} and
            \cref{dfn:set-builder-notation}}\\
            &\quad\text{De Morgan's Law}\\
            &\quad\text{Disjunctive Syllogism}\\
            &\quad\text{Conjunction Introduction}\\
            &\quad\text{De Morgan's Law}\\
            &\quad\text{Existential Generalization}\\
            &\quad\text{Logical Equivalence}\\
        \end{aligned}
    \end{equation*}
    By \cref{dfn:subset}, the last statement is equivalent to \( \powerset{A} \nsubseteq A \).
\end{proof}

\begin{thm}
    \label{thm:power-set-of-A-is-never-equal-to-A}
    For any set \( A \), \( \powerset{A} \neq A \).
\end{thm}

\begin{proof}
    Let \( A \) be an arbitrary set. It follows from
    \cref{thm:power-set-of-A-can-never-be-subset-of-A}
    that \( \powerset{A} \nsubseteq A \). Thus,
    \( \powerset{A} \neq A \) by \cref{thm:set-equality-in-terms-of-subsets}.
\end{proof}

\begin{thm}
    \label{thm:A-is-a-subset-of-B-iff-the-power-set-of-A-is-the-subset-of-the-power-set-of-B}
    For any two sets \( A \) and \( B \), \( A \subseteq B \iff \powerset{A} \subseteq \powerset{B} \).
    \begin{equation*}
        \forall A,B \, \Big[ A \subseteq B \iff \powerset{A} \subseteq \powerset{B} \Big]
    \end{equation*}
\end{thm}

\begin{proof}
    Let \( A \) and \( B \) be arbitrary sets. Suppose that \( A \subseteq B \). By \cref{dfn:subset}, it
    follows that,
    \begin{equation}
        \label{eq:state-that-A-is-a-subset-of-B-in-logical-form}
        \forall x \, \mleft( x \in A \implies x \in B \mright)
    \end{equation}
    Let \( z_{0} \in \powerset{A} \) be arbitrary. By \cref{dfn:power-set,dfn:subset}, it follows that
    \( z_{0} \subseteq A \). Again, by \cref{dfn:subset}, we have,
    \begin{equation}
        \label{eq:state-that-z0-is-a-subset-of-A-in-logical-form}
        \forall x \, \mleft( x \in z_{0} \implies x \in A \mright)
    \end{equation}
    Now, let \( \alpha \in z_{0} \) be arbitrary. By \eqref{eq:state-that-z0-is-a-subset-of-A-in-logical-form},
    it follows that \( \alpha \in A \) and thus \( \alpha \in B \) by \eqref{eq:state-that-A-is-a-subset-of-B-in-logical-form}.
    \begin{equation*}
        \therefore \, \alpha \in z_{0} \implies \alpha \in B
    \end{equation*}
    Since \( \alpha \) is arbitrary, we have \( \forall x \, \mleft( x \in z_{0} \implies x \in B \mright) \) by Universal Generalization.
    In other words, \( z_{0} \subseteq B \) and \( z_{0} \in \powerset{B} \).
    \begin{equation*}
        \therefore \, z_{0} \in \powerset{A} \implies z_{0} \in \powerset{B}
    \end{equation*}
    Since \( z_{0} \) is arbitrary, we can generalize again,
    \begin{equation*}
        \forall x \, \mleft( x \in \powerset{A} \implies x \in \powerset{B} \mright)
    \end{equation*}
    By \cref{dfn:subset}, it follows that \( \powerset{A} \subseteq \powerset{B} \).
    \begin{equation}
        \label{eq:A-subset-B-implies-powerset-A-subset-powerset-B}
        \therefore \, A \subseteq B \implies \powerset{A} \subseteq \powerset{B}
    \end{equation}
    To prove the converse, suppose that \( \powerset{A} \subseteq \powerset{B} \). By \cref{thm:every-set-is-an-element-of-its-powerset},
    it follows that \( A \in \powerset{A} \). Since \( \powerset{A} \subseteq \powerset{B} \), this means that \( A \in \powerset{B} \),
    which in turn implies \( A \subseteq B \).
    \begin{equation}
        \label{eq:powerset-A-subset-powerset-B-implies-A-subset-B}
        \therefore \, \powerset{A} \subseteq \powerset{B} \implies A \subseteq B
    \end{equation}
    It follows from \eqref{eq:A-subset-B-implies-powerset-A-subset-powerset-B} and \eqref{eq:powerset-A-subset-powerset-B-implies-A-subset-B}
    that \( A \subseteq B \iff \powerset{A} \subseteq \powerset{B} \). Since \( A \) and \( B \) are arbitrary, the theorem follows from
    Universal Generalization.
\end{proof}

\begin{thm}[Transitivity]
    \label{thm:transitivity-of-the-subset-relation}
    For any sets \( A \), \( B \), and \( C \), if
    \( A \subseteq B \) and \( B \subseteq C \), then
    \( A \subseteq C \).
    \begin{equation*}
        \forall A,B,C \, \mleft( A \subseteq B \land
        B \subseteq C \implies A \subseteq C \mright)
    \end{equation*}
\end{thm}

\begin{proof}
    Let \( A \), \( B \), and \( C \) be arbitrary sets such that
    \( A \subseteq B \) and \( B \subseteq C \). The following chain
    of inferences completes the proof:
    \begin{equation*}
        \begin{aligned}
            A \subseteq B \land B \subseteq C& \\
            A \subseteq B& \\
            B \subseteq C& \\
            \forall x \, \mleft( x \in A \implies x \in B \mright)& \\
            \forall x \, \mleft( x \in B \implies x \in C \mright)& \\
            \alpha \in A \implies \alpha \in B& \\
            \alpha \in B \implies \alpha \in C& \\
            \alpha \in A \implies \alpha \in C& \\
            \forall x \, \mleft( x \in A \implies x \in C \mright)&\\
            A \subseteq C& \\
            \therefore \, A \subseteq B \land B \subseteq C
            \implies A \subseteq C& \\
        \end{aligned}
        \text{  }
        \begin{aligned}
            &\quad\text{Assumption} \\
            &\quad\text{Conjunction Elimination} \\
            &\quad\text{Conjunction Elimination} \\
            &\quad\text{\cref{dfn:subset}} \\
            &\quad\text{\cref{dfn:subset}} \\
            &\quad\text{Universal Instantiation} \\
            &\quad\text{Universal Instantiation} \\
            &\quad\text{Hypothetical Syllogism} \\
            &\quad\text{Universal Generalization} \\
            &\quad\text{\cref{dfn:subset}} \\
            &\quad\text{Deduction Theorem} \\
        \end{aligned}
    \end{equation*}
    Since \( A \), \( B \), and \( C \) are arbitrary, the theorem follows
    by Universal Generalization.
\end{proof}

\begin{thm}
    \label{thm:every-member-of-a-system-of-sets-is-a-subset-of-the-union}
    For any system of sets \( S \), if \( A \in S \) then, \( A \subseteq \Union S \).
    \begin{equation*}
        \forall S \, \forall A \, \mleft( A \in S \implies A \subseteq \Union S \mright)
    \end{equation*}
\end{thm}

\begin{proof}
    Let \( S \) be an arbitrary system of sets. Let \( A \) be a set such that \( A \in S \).
    Let \( x_{0} \) be an object such that \( x_{0} \in A \).
    \begin{equation*}
        \begin{aligned}
            x_{0} \in A& \\
            A \in S& \\
            \exists \, T \,{\in}\, S \, \mleft( x_{0} \in T \mright)& \\
            x_{0} \in \union S& \\
            x_{0} \in A \implies x_{0} \in \union S& \\
            \forall x \, \mleft( x \in A \implies x \in \union S \mright)& \\
            A \subseteq \union S& \\
            \therefore \, A \in S \implies A \subseteq \union S& \\
        \end{aligned}
        \text{  }
        \begin{aligned}
            &\quad\text{Assumption}\\
            &\quad\text{Assumption} \\
            &\quad\text{Existential Generalization}\\
            &\quad\text{\cref{dfn:union-of-sets}}\\
            &\quad\text{Deduction Theorem}\\
            &\quad\text{Universal Generalization}\\
            &\quad\text{\cref{dfn:subset}}\\
            &\quad\text{Deduction Theorem}
        \end{aligned}
    \end{equation*}
    Since \( A \) and \( S \) are arbitrary, the theorem follows from Universal
    Generalization.
\end{proof}

\begin{thm}
    \label{thm:a-set-is-a-subset-of-its-union-with-another-set}
    For any two sets \( A \) and \( B \), \( A \subseteq A \union B \)
    and \( B \subseteq A \union B \).
    \begin{equation*}
        \forall A,B \, \mleft( A \subseteq A \union B \land B \subseteq A \union B \mright)
    \end{equation*}
\end{thm}

\begin{proof}
    This is just a special case of \cref{thm:every-member-of-a-system-of-sets-is-a-subset-of-the-union}
    with \( S = \mleft\{ A,B \mright\} \). By \cref{dfn:unordered-pair}, we have \( A \in \mleft\{ A,B \mright\}
    \land B \in \mleft\{ A,B \mright\} \).
    By \cref{thm:every-member-of-a-system-of-sets-is-a-subset-of-the-union},
    \begin{gather*}
        A \subseteq \Union \mleft\{ A,B \mright\} = A \union B \\
        B \subseteq \Union \mleft\{ A,B \mright\} = A \union B
    \end{gather*}
    Since \( A \) and \( B \) are arbitrary, the theorem follows from Universal Generalization.
\end{proof}

\begin{thm}
    \label{thm:the-intersection-of-any-system-of-sets-is-a-subset-of-each-member-of-the-system}
    For any non-empty system of sets \( S \),
    \begin{equation*}
        \forall A \, \mleft( A \in S \implies \Intersect S \subseteq A \mright)
    \end{equation*}
\end{thm}

\begin{proof}
    Let \( S \) be a system of sets such that \( S \neq \emptyset \).
    \begin{equation*}
        \begin{aligned}
            A \in S& \\
            \forall z \, \mleft( z \in S \implies x_{0} \in z \mright)&\\
            A \in S \implies x_{0} \in A& \\
            x_{0} \in A& \\
            \forall z \, \mleft( z \in S \implies x_{0} \in z \mright) \implies
            x_{0} \in A& \\
            \intersect S \subseteq A& \\
            \therefore \, A \in S \implies \intersect S \subseteq A& \\
        \end{aligned}
        \text{  }
        \begin{aligned}
            &\quad\text{Assumption}\\
            &\quad\text{Assumption, \cref{dfn:intersection-of-sets}}\\
            &\quad\text{Universal Instantiation}\\
            &\quad\text{Modus Ponens}\\
            &\quad\text{Deduction Theorem}\\
            &\quad\text{Universal Generalization, \cref{dfn:intersection-of-sets}}\\
            &\quad\text{Deduction Theorem}
        \end{aligned}
    \end{equation*}
    Since \( A \) and \( S \) are arbitrary, the theorem follows from Universal Generalization.
\end{proof}

\begin{thm}
    \label{thm:the-intersection-of-two-sets-is-a-subset-of-each-set}
    For any two sets \( A \) and \( B \), \( A \intersect B \subseteq A \) and
    \( A \intersect B \subseteq B \).
    \begin{equation*}
        \forall A,B \, \mleft( A \intersect B \subseteq A \land A \intersect B \subseteq B \mright)
    \end{equation*}
\end{thm}

\begin{proof}
    This is a special case of \cref{thm:the-intersection-of-any-system-of-sets-is-a-subset-of-each-member-of-the-system}
    with \( S = \mleft\{ A,B \mright\} \). By \cref{dfn:unordered-pair}, it follows that \( A \in \mleft\{ A,B \mright\}
    \land B \in \mleft\{ A,B \mright\} \). Thus, \( A \intersect B = \Intersect \mleft\{ A,B \mright\} \subseteq A \) and
    \( A \intersect B = \Intersect \mleft\{ A,B \mright\} \subseteq B \) by \cref{thm:the-intersection-of-any-system-of-sets-is-a-subset-of-each-member-of-the-system}
    as required. Universal Generalization completes the proof.
\end{proof}

\begin{thm}
    \label{thm:subset-in-terms-of-union}
    For any two sets \( A \) and \( B \), \( A \subseteq B \) if and only if
    \( A \union B = B \).
    \begin{equation*}
        \forall A,B \, \mleft( A \subseteq B \iff A \union B = B \mright)
    \end{equation*}
\end{thm}

\begin{proof}
    Let \( A \) and \( B \) be arbitrary sets.
    \begin{equation*}
        \begin{aligned}
            A \subseteq B& \\
            \forall x \, \mleft( x \in A \implies x \in B \mright)& \\
            x_{0} \in A \union B& \\
            x_{0} \in A \lor x_{0} \in B& \\
            x_{0} \in A \implies x_{0} \in B& \\
            x_{0} \in B \implies x_{0} \in B& \\
            x_{0} \in B& \\
            x_{0} \in A \union B \implies x_{0} \in B& \\
            \forall x \, \mleft( x \in A \union B \implies x \in B \mright)& \\
            A \union B \subseteq B& \\
            \therefore \, A \subseteq B \implies A \union B \subseteq B&
        \end{aligned}
        \text{  }
        \begin{aligned}
            &\quad\text{Assumption}\\
            &\quad\text{\cref{dfn:subset}}\\
            &\quad\text{Assumption}\\
            &\quad\text{\cref{thm:union-in-terms-of-logical-or}}\\
            &\quad\text{Universal Instantiation}\\
            &\quad\text{Tautology}\\
            &\quad\text{Disjunction Elimination}\\
            &\quad\text{Deduction Theorem}\\
            &\quad\text{Universal Generalization}\\
            &\quad\text{\cref{dfn:subset}}\\
            &\quad\text{Deduction Theorem}
        \end{aligned}
    \end{equation*}
    By \cref{thm:a-set-is-a-subset-of-its-union-with-another-set}, we know that
    \( B \subseteq A \union B \) is also true.
    \begin{equation*}
        \therefore \, A \subseteq B \implies A \union B \subseteq B \land
        B \subseteq A \union B
    \end{equation*}
    By \cref{thm:set-equality-in-terms-of-subsets}, we have,
    \begin{equation}
        \label{eq:A-subset-B-implies-A-union-B-equals-B}
        A \subseteq B \implies A \union B = B
    \end{equation}
    To prove the converse, suppose that \( A \union B = B \).
    \begin{equation*}
        \begin{aligned}
            A \union B = B& \\
            \forall x \, \mleft( x \in A \lor x \in B \iff x \in B \mright)& \\
            x_{0} \in A& \\
            x_{0} \in A \lor x_{0} \in B \iff x_{0} \in B& \\
            x_{0} \in A \lor x_{0} \in B& \\
            x_{0} \in B& \\
            x_{0} \in A \implies x_{0} \in B& \\
            \forall x \, \mleft( x \in A \implies x \in B \mright)& \\
            \therefore \, A \subseteq B& \\
        \end{aligned}
        \text{  }
        \begin{aligned}
            &\quad\text{Assumption}\\
            &\quad\text{\nameref{axm:axiom-of-extensionality}, \cref{thm:union-in-terms-of-logical-or}}\\
            &\quad\text{Assumption}\\
            &\quad\text{Universal Instantiation}\\
            &\quad\text{Disjunction Introduction}\\
            &\quad\text{Modus Ponens}\\
            &\quad\text{Deduction Theorem}\\
            &\quad\text{Universal Generalization}\\
            &\quad\text{\cref{dfn:subset}}\\
        \end{aligned}
    \end{equation*}
    By the Deduction Theorem,
    \begin{equation}
        \label{eq:A-union-B-equals-B-implies-A-subset-B}
        A \union B = B \implies A \subseteq B
    \end{equation}
    Together, \eqref{eq:A-subset-B-implies-A-union-B-equals-B} and \eqref{eq:A-union-B-equals-B-implies-A-subset-B}
    establishes the biconditional \( A \subseteq B \iff A \union B = A \). Since \( A \) and \( B \) are arbitrary,
    the theorem follows from Universal Generalization.
\end{proof}

\begin{thm}
    \label{thm:subset-in-terms-of-intersection}
    For any two sets \( A \) and \( B \), \( A \subseteq B \) if and only if
    \( A \intersect B = A \).
    \begin{equation*}
        \forall A,B \, \mleft( A \subseteq B \iff A \intersect B = A \mright)
    \end{equation*}
\end{thm}

\begin{proof}
    Let \( A \) and \( B \) be arbitrary sets. \cref{thm:the-intersection-of-two-sets-is-a-subset-of-each-set}
    dictates that,
    \begin{equation*}
        A \intersect B \subseteq A
    \end{equation*}
    It remains to prove the inclusion relation in reverse.
    \begin{equation*}
        \begin{aligned}
            A \subseteq B& \\
            \forall x \, \mleft( x \in A \implies x \in B \mright)& \\
            x_{0} \in A& \\
            x_{0} \in A \implies x_{0} \in B& \\
            x_{0} \in B& \\
            x_{0} \in A \land x_{0} \in B& \\
            x_{0} \in A \intersect B& \\
            x_{0} \in A \implies x_{0} \in A \intersect B& \\
            \forall x \, \mleft( x \in A \implies x \in A \intersect B \mright)& \\
            \therefore \, A \subseteq A \intersect B& \\
        \end{aligned}
        \text{  }
        \begin{aligned}
            &\quad\text{Assumption}\\
            &\quad\text{\cref{dfn:subset}}\\
            &\quad\text{Assumption}\\
            &\quad\text{Universal Instantiation}\\
            &\quad\text{Modus Ponens}\\
            &\quad\text{Conjunction Introduction}\\
            &\quad\text{\cref{thm:intersection-in-terms-of-logical-and}}\\
            &\quad\text{Deduction Theorem}\\
            &\quad\text{Universal Generalization}\\
            &\quad\text{\cref{dfn:subset}}\\
        \end{aligned}
    \end{equation*}
    Since \( A \intersect B \subseteq A \land A \subseteq A \intersect B \),
    it follows from \cref{thm:set-equality-in-terms-of-subsets} that
    \( A \intersect B = A \).
    \begin{equation}
        \label{eq:A-subset-B-implies-A-intersect-B-equals-A}
        \therefore \, A \subseteq B \implies A \intersect B = A
    \end{equation}
    The proof for the converse is fairly straightforward:
    \begin{equation*}
        \begin{aligned}
            A \intersect B = A& \\
            \forall x \, \mleft( x \in A \land x \in B \iff x \in A \mright)& \\
            x_{0} \in A& \\
            x_{0} \in A \land x_{0} \in B \iff x_{0} \in A& \\
            x_{0} \in A \land x_{0} \in B& \\
            x_{0} \in B& \\
            x_{0} \in A \implies x_{0} \in B& \\
            \forall x \, \mleft( x \in A \implies x \in B \mright)& \\
            A \subseteq B&\\
            A \intersect B = A \implies A \subseteq B& \\
        \end{aligned}
        \text{  }
        \begin{aligned}
            &\quad\text{Assumption}\\
            &\quad\text{\nameref{axm:axiom-of-extensionality}, \cref{thm:intersection-in-terms-of-logical-and}}\\
            &\quad\text{Assumption}\\
            &\quad\text{Universal Instantiation}\\
            &\quad\text{Modus Ponens}\\
            &\quad\text{Conjunction Elimination}\\
            &\quad\text{Deduction Theorem}\\
            &\quad\text{Universal Generalization}\\
            &\quad\text{\cref{dfn:subset}}\\
            &\quad\text{Deduction Theorem}
        \end{aligned}
    \end{equation*}
    This along with \eqref{eq:A-subset-B-implies-A-intersect-B-equals-A} establish the biconditional,
    \begin{equation*}
        A \subseteq B \iff A \intersect B = A
    \end{equation*}
    Since \( A \) and \( B \) are arbitrary, the theorem follows from Universal Generalization.
\end{proof}

\begin{thm}
    \label{set-identities:subset-in-terms-of-empty-difference}
    For any two sets \( A \) and \( B \), \( A \subseteq B \iff A - B = \emptyset\).
    \begin{equation*}
        \forall A,B \, \mleft( A \subseteq B \iff A - B = \emptyset \mright)
    \end{equation*}
\end{thm}

\begin{proof}
    Let \( A \) and \( B \) be arbitrary sets. The theorem can be derived by applying the laws
    of propositional logic to \cref{dfn:subset}:
    \begin{align*}
        A \subseteq B &\iff \forall x \, \mleft( x \in A \implies x \in B \mright)\\
        &\iff \forall x \, \mleft( \lnot \, \mleft( x \in A \land x \notin B \mright) \mright) \\
        &\iff \forall x \, \mleft( x \notin A - B \mright) \\
        &\iff A - B = \emptyset \qedhere
    \end{align*}
\end{proof}

\begin{thm}
    \label{set-identities:left-distributivity-of-subset-relation-over-intersection}
    For any three sets \( A \), \( B \), and \( C \), \( A \subseteq \mleft( B \intersect C \mright)
    \iff A \subseteq B \land A \subseteq C \).
    \begin{equation*}
        \forall A,B,C \, \Big( A \subseteq \mleft( B \intersect C \mright) \iff
        A \subseteq B \land A \subseteq C \Big)
    \end{equation*}
\end{thm}

\begin{proof}
    Let \( A \), \( B \), and \( C \) be arbitrary sets. Using elementary propositional logic, \cref{dfn:subset,thm:intersection-in-terms-of-logical-and},
    we can re-express \( A \subseteq \mleft( B \intersect C \mright) \) in another logically equivalent form:
    \begin{align*}
        A \subseteq \mleft( B \intersect C \mright)
        &\iff
        \forall x \, \mleft( x \in A \implies x \in B \land x \in C \mright)\\
        &\iff
        \forall x \, \mleft( x \notin A \,\,{\lor} \mleft( x \in B \land x \in C \mright) \mright)\\
        &\iff
        \forall x \, \mleft( \mleft( x \notin A \lor x \in B \mright)
        \land \mleft( x \notin A \lor x \in C \mright) \mright)\\
        &\iff
        \forall x \, \mleft( x \notin A \lor x \in B \mright)
        \land
        \forall x \, \mleft( x \notin A \lor x \in C \mright)\\
        &\iff
        \forall x \, \mleft( x \in A \implies x \in B \mright)
        \land
        \forall x \, \mleft( x \in A \implies x \in C \mright)\\
        &\iff
        A \subseteq B \land A \subseteq C
    \end{align*}
    Since \( A \), \( B \), and \( C \) are arbitrary, the theorem follows from Universal Generalization.
\end{proof}

\begin{thm}
    \label{set-identities:right-distributivity-of-subset-relation-over-union}
    For any three sets \( A \), \( B \), and \( C \), \( \mleft( A \union B \mright) \subseteq C
    \iff A \subseteq C \land B \subseteq C\).
    \begin{equation*}
        \forall A,B,C \, \Big( \mleft( A \union B \mright) \subseteq C \iff
        A \subseteq C \land B \subseteq C \Big)
    \end{equation*}
\end{thm}

\begin{proof}
    Let \( A \), \( B \), and \( C \) be arbitrary sets. Using elementary propositional logic, \cref{dfn:subset,thm:union-in-terms-of-logical-or},
    we can re-express \( \mleft( A \union B \mright) \subseteq C \) in another logically equivalent form:
    \begin{align*}
        \mleft( A \union B \mright) \subseteq C
        &\iff
        \forall x \, \big( \mleft( x \in A \lor x \in B \mright)
        \implies x \in C \big)\\
        &\iff
        \forall x \, \big( \mleft( x \notin A \land x \notin B \mright) \lor x \in C \big)\\
        &\iff
        \forall x \, \big( \mleft( x \notin A \lor x \in C \mright)
        \land
        \mleft( x \notin B \lor x \in C \mright) \big)\\
        &\iff
        \forall x \, \big( x \in A \implies x \in C \land x \in B \implies x \in C \big)\\
        &\iff
        \forall x \, \mleft( x \in A \implies x \in C \mright) \land
        \forall x \, \mleft( x \in B \implies x \in C  \mright)\\
        &\iff
        A \subseteq C \land B \subseteq C
    \end{align*}
    Since \( A \), \( B \), and \( C \) are arbitrary, the theorem follows from Universal Generalization.
\end{proof}

\begin{thm}
    \label{thm:the-union-of-subsets-is-a-subset-of-the-union}
    Let \( A \), \( B \), \( C \), and \( D \) be arbitrary sets. If \( A \subseteq C \) and \( B \subseteq D \), then
    \( A \union B \subseteq C \union D \).
    \begin{equation*}
        \forall A,B,C,D \, \Big( \mleft( A \subseteq C \mright) \land \mleft( B \subseteq D \mright) \implies
        \mleft( A \union B \mright) \subseteq \mleft( C \union D \mright) \Big)
    \end{equation*}
\end{thm}

\begin{proof}
    Let \( x_{0} \) be an arbitrary object. Since \( A \subseteq C \) and \( B \subseteq D \), it follows from \cref{dfn:subset}
    that,
    \begin{gather*}
        x_{0} \in A \implies x_{0} \in C \\
        x_{0} \in B \implies x_{0} \in D
    \end{gather*}
    Suppose that \( x_{0} \in A \union B \). Then by \cref{thm:union-in-terms-of-logical-or},
    \begin{equation*}
        x_{0} \in A \lor x_{0} \in B
    \end{equation*}
    By Constructive Dilemma, it follows that \( x_{0} \in C \lor x_{0} \in D \).
    \begin{equation*}
        \therefore \, x_{0} \in \mleft( A \union B \mright) \implies x_{0} \in \mleft( C \union D \mright)
    \end{equation*}
    The theorem follows by Universal Generalization
\end{proof}

\subsection{Union}
\begin{thm}
    \label{set-identities:union-of-empty-set-is-empty}
    \( \Union \emptyset = \emptyset \).
\end{thm}

\begin{proof}
    Let \( x \) be an arbitrary object and \( T \) be any arbitrary set. Thus, we can make the following inferences:
    \begin{equation*}
        \begin{aligned}[c]
            T \notin \emptyset& \\
            T \notin \emptyset \lor x \notin T& \\
            \lnot \, \mleft( T \in \emptyset \land x \in T \mright)& \\
            \forall T \, \big[ \lnot \, \mleft( T \in \emptyset \land x \in T \mright) \big]& \\
            \lnot \, \exists \, T \big[ T \in \emptyset \land x \in T \big]& \\
        \end{aligned}
        \text{  }
        \begin{aligned}[c]
            &\quad\text{\cref{dfn:empty-set}} \\
            &\quad\text{Disjunction Introduction} \\
            &\quad\text{De Morgan's Law} \\
            &\quad\text{Universal Generalization} \\
            &\quad\text{Logical Equivalence}
        \end{aligned}
    \end{equation*}
    By \cref{dfn:union-of-sets}, the last statement implies \( x \notin \Union \emptyset \).
    Since \( x \) is arbitrary, we conclude by Universal Generalization that,
    \begin{equation*}
        \forall x \, \mleft( x \notin \Union \emptyset \mright)
    \end{equation*}
    \( \Union \emptyset = \emptyset \) follows readily from \cref{dfn:empty-set}.
\end{proof}

\begin{thm}
    \label{thm:union-of-a-singleton}
    For any set \( A \), \( \Union \mleft\{ A \mright\} = A \)
\end{thm}

\begin{proof}
    Let \( A \) be an arbitrary set. By \cref{dfn:union-of-sets},
    \begin{equation}
        \label{eq:express-the-union-of-singleton-using-first-order-logic}
        \forall x \, \Big[ x \in \Union \mleft\{ A \mright\}
        \iff
        \exists S \, \big( S \in \mleft\{ A \mright\} \land x \in S \big) \Big]
    \end{equation}
    But \( \forall x \, \big( x \in \mleft\{ A \mright\} \iff x = A \big) \) by \cref{dfn:unordered-pair}.
    Thus, \eqref{eq:express-the-union-of-singleton-using-first-order-logic} can be equivalently re-expressed as,
    \begin{equation*}
        \forall x \, \Big[ x \in \Union \mleft\{ A \mright\}
        \iff
        \exists S \, \big( S = A \land x \in S \big)
        \Big]
    \end{equation*}
    By \cref{logic:using-existential-quantifier-to-express-the-fact-that-an-object-satisfies-a-predicate},
    this becomes,
    \begin{equation*}
        \forall x \, \Big[ x \in \Union \mleft\{ A \mright\}
        \iff
        x \in A \Big]
    \end{equation*}
    By the \nameref{axm:axiom-of-extensionality}, we thus have,
    \begin{equation*}
        \Union \mleft\{ A \mright\} = A \qedhere
    \end{equation*}
\end{proof}

\begin{thm}[Idempotence]
    \label{set-identities:idempotence-of-unions}
    For any set \( A \), \( A \union A = A\).
    \begin{equation*}
        \forall A \, \mleft( A \union A = A \mright)
    \end{equation*}
\end{thm}

\begin{proof}
    Let \( A \) be an arbitrary set. Let \( x_{0} \) be arbitrary. By \cref{thm:union-in-terms-of-logical-or},
    \begin{align*}
        x_{0} \in A \union A &\iff x_{0} \in A \lor x_{0} \in A\\
        &\iff x_{0} \in A
    \end{align*}
    Since \( x_{0} \) is arbitrary, we conclude by Universal Generalization that,
    \begin{equation*}
        \forall x \, \mleft( x \in A \union A \iff x \in A \mright)
    \end{equation*}
    By \nameref{axm:axiom-of-extensionality}, we have \( A \union A = A \).
\end{proof}

\begin{thm}[Commutativity]
    \label{set-identities:commutativity-of-union}
    For any two sets \( A \) and \( B \), \( A \union B = B \union A \).
    \begin{equation*}
        \forall A,B \, \mleft( A \union B = B \union A \mright)
    \end{equation*}
\end{thm}

\begin{proof}
    Let \( A \) and \( B \) be two arbitrary sets. Let \( x_{0} \) be an arbitrary object.
    The logical \( \lor \)-operator is commutative, thus,
    \begin{equation*}
        x_{0} \in A \lor x_{0} \in B \iff x_{0} \in B \lor x_{0} \in A
    \end{equation*}
    Since \( x_{0} \) is arbitrary, Universal Generalization permits us to infer,
    \begin{equation}
        \label{eq:quantified-biconditional-for-commutativity-of-logical-or}
        \forall x \, \mleft( x \in A \lor x \in B \iff x \in B \lor x \in A \mright)
    \end{equation}
    Applying \cref{thm:union-in-terms-of-logical-or} to \eqref{eq:quantified-biconditional-for-commutativity-of-logical-or}
    yields,
    \begin{equation*}
        \forall x \, \mleft( x \in A \union B \iff x \in B \union A \mright)
    \end{equation*}
    And finally, by the \nameref{axm:axiom-of-extensionality}, \( A \union B = B \union A \).
\end{proof}

\begin{thm}
    \label{set-identities:any-set-union-empty-set-equals-the-set-itself}
    For any set \( A \), \( A \union \emptyset = \emptyset \union A = A \)
\end{thm}

\begin{proof}
    Let \( A \) be an arbitrary set. \cref{set-identities:commutativity-of-union} already implies \( A \union \emptyset = \emptyset \union A \).
    Thus, it suffices to prove that \( A \union \emptyset = A \). Suppose that \( x_{0} \in A \union \emptyset \) for some arbitrary \( x_{0} \). By
    \cref{thm:union-in-terms-of-logical-or}, we have \( x_{0} \in A \lor x_{0} \in \emptyset\). Since \( \forall x \, \mleft( x \notin \emptyset \mright) \),
    we conclude \( x_{0} \in A \) by Disjunction Elimination. Thus, \( x_{0} \in A \union \emptyset \implies x_{0} \in A \). To prove the converse, suppose that
    \( x_{0} \in A \); then \( x_{0} \in A \lor x_{0} \in \emptyset \) is obviously true. Thus, \( x_{0} \in A \implies x_{0} \in A \union \emptyset \).
    \begin{equation*}
        \therefore \, x_{0} \in A \union \emptyset \iff x_{0} \in A
    \end{equation*}
    The theorem follows from Universal Generalization and \nameref{axm:axiom-of-extensionality}.
\end{proof}

\begin{thm}[Associativity]
    \label{set-identities:associativity-of-unions}
    \( A \union \mleft( B \union C \mright) = \mleft( A \union B \mright) \union C \) for any arbitrary sets \( A \), \( B \), and \( C \).
    \begin{equation*}
        \forall A,B,C \, \big( A \union \mleft( B \union C \mright) = \mleft( A \union B \mright) \union C \big)
    \end{equation*}
\end{thm}

\begin{proof}
    Let \( A \), \( B \), and \( C \) be arbitrary sets. Let \( x_{0} \) be an arbitrary object. By the laws of propositional logic,
    \begin{equation}
        \label{eq:associativity-of-logical-or-operator}
        x_{0} \in A \lor \mleft( x_{0} \in B \lor x_{0} \in C \mright)
        \iff
        \mleft( x_{0} \in A \lor x_{0} \in B \mright) \lor x_{0} \in C
    \end{equation}
    Applying \cref{thm:union-in-terms-of-logical-or} to each formula in the parentheses, we have,
    \begin{align}
        x_{0} \in B \union C &\iff x_{0} \in B \lor x_{0} \in C
        \label{eq:B-union-C-in-terms-of-B-or-C} \\
        x_{0} \in A \union B &\iff x_{0} \in A \lor x_{0} \in B
        \label{eq:A-union-B-in-terms-of-A-or-B}
    \end{align}
    Applying \eqref{eq:B-union-C-in-terms-of-B-or-C} and \eqref{eq:A-union-B-in-terms-of-A-or-B}
    into \eqref{eq:associativity-of-logical-or-operator} yields,
    \begin{align}
        x_{0} \in A \lor x_{0} \in B \union C &\iff
        x_{0} \in A \,\,{\lor} \mleft( x_{0} \in B \lor x_{0} \in C \mright)
        \label{eq:A-union-B-union-C-in-terms-of-logical-or}\\
        x_{0} \in A \union B \lor x_{0} \in C &\iff
        \mleft( x_{0} \in A \lor x_{0} \in B \mright) \lor x_{0} \in C
        \label{eq:A-union-B-group-union-C-in-terms-of-logical-or}
    \end{align}
    Finally, applying \cref{thm:union-in-terms-of-logical-or} to \eqref{eq:A-union-B-union-C-in-terms-of-logical-or}
    and \eqref{eq:A-union-B-group-union-C-in-terms-of-logical-or} yields,
    \begin{align*}
        x_{0} \in A \union \mleft( B \union C \mright)
        &\iff
        x_{0} \in A \lor x_{0} \in B \union C\\
        &\iff
        x_{0} \in A \,\,{\lor} \mleft( x_{0} \in B \lor x_{0} \in C \mright)\\
        &\iff
        \mleft( x_{0} \in A \lor x_{0} \in B \mright) \lor x_{0} \in C\\
        &\iff
        x_{0} \in A \union B \lor x_{0} \in C\\
        &\iff
        x_{0} \in \mleft( A \union B \mright) \union C
    \end{align*}
    Since \( x_{0} \) is arbitrary, Universal Generalization yields,
    \begin{equation*}
        \forall x \, \Big( x \in A \union \mleft( B \union C \mright) \iff
        x \in \mleft( A \union B \mright) \union C  \Big)
    \end{equation*}
    Thus, \( A \union \mleft( B \union C \mright) = A \union \mleft( B \union C \mright) \) by the \nameref{axm:axiom-of-extensionality}.
    Since \( A \), \( B \), and \( C \) are arbitrary, the theorem follows from Universal Generalization.
\end{proof}
\subsection{Intersection}
\begin{thm}[Idempotence]
    \label{set-identities:idempotence-of-intersections}
    For any set \( A \), \( A \intersect A = A \)
    \begin{equation*}
        \forall A \, \mleft( A \intersect A = A \mright)
    \end{equation*}
\end{thm}

\begin{proof}
    Let \( A \) be an arbitrary set. Let \( x_{0} \) be an arbitrary object.
    By elementary propositional logic and \cref{thm:intersection-in-terms-of-logical-and},
    \begin{align*}
        x_{0} \in A
        &\iff
        x_{0} \in A \land x_{0} \in A \\
        &\iff
        x_{0} \in A \intersect A
    \end{align*}
    By Universal Generalization and \nameref{axm:axiom-of-extensionality}, \( A = A \intersect A \).
\end{proof}

\begin{thm}[Commutativity]
    \label{set-identities:commutativity-of-intersection}
    For any two sets \( A \) and \( B \), \( A \intersect B = B \intersect A\).
    \begin{equation*}
        \forall A,B \, \mleft( A \intersect B = B \intersect A \mright)
    \end{equation*}
\end{thm}

\begin{proof}
    Let \( A \) and \( B \) be two arbitrary sets. Let \( x_{0} \) be an arbitrary object.
    By elementary propositional logic and \cref{thm:intersection-in-terms-of-logical-and},
    \begin{align*}
        x_{0} \in A \intersect B
        &\iff
        x_{0} \in A \land x_{0} \in B \\
        &\iff
        x_{0} \in B \land x_{0} \in A \\
        &\iff
        x_{0} \in B \intersect A
    \end{align*}
    Since \( x_{0} \) is arbitrary, we conclude by Universal Generalization that,
    \begin{equation*}
        \forall x \, \mleft( x \in A \intersect B \iff x \in B \intersect A \mright)
    \end{equation*}
    Finally, the \nameref{axm:axiom-of-extensionality} implies \( A \intersect B = B \intersect A \).
    Since \( A \) and \( B \) are arbitrary, the theorem follows from Universal Generalization.
\end{proof}

\begin{thm}[Absorption]
    \label{set-identities:intersection-of-any-set-with-empty-set-is-always-empty}
    For any set \( A \), \( A \intersect \emptyset = \emptyset \intersect A = \emptyset \).
    \begin{equation*}
        \forall A \, \mleft( A \intersect \emptyset = \emptyset \intersect A = \emptyset \mright)
    \end{equation*}
\end{thm}

\begin{proof}
    Let \( A \) be an arbitrary set. \cref{set-identities:commutativity-of-intersection} already implies
    \( A \intersect \emptyset = \emptyset \intersect A \).
    Thus, it suffices to prove that \( A \intersect \emptyset = \emptyset \). Let \( x_{0} \) be
    an arbitrary object. \cref{thm:intersection-in-terms-of-logical-and} implies that,
    \begin{equation*}
        x_{0} \in A \intersect \emptyset \iff x_{0} \in A \land x_{0} \in \emptyset
    \end{equation*}
    But \cref{dfn:empty-set} implies that \( \forall x \, \mleft( x \notin \emptyset \mright) \).
    \begin{equation*}
        \therefore \, x_{0} \notin \emptyset
    \end{equation*}
    In other words, \( \lnot \, \mleft( x_{0} \in A \land x_{0} \in \emptyset \mright) \).
    \begin{equation*}
        \therefore \, x_{0} \notin A \intersect \emptyset
    \end{equation*}
    Since \( x_{0} \) is arbitrary, we conclude by Universal Generalization that,
    \begin{equation*}
        \forall x \, \mleft( x \notin A \intersect \emptyset \mright)
    \end{equation*}
    By \cref{dfn:empty-set}, \( A \intersect \emptyset = \emptyset \).
\end{proof}

\begin{thm}[Associativity]
    \label{set-identities:associativity-of-intersection}
    For any three sets \( A \), \( B \), and \( C \), \( A \intersect \mleft( B \intersect C \mright) = \mleft( A \intersect B \mright) \intersect C \).
    \begin{equation*}
        \forall A,B,C \, \Big( A \intersect \mleft( B \intersect C \mright) = \mleft( A \intersect B \mright) \intersect C \Big)
    \end{equation*}
\end{thm}

\begin{proof}
    Let \( A \), \( B \), and \( C \) be arbitrary sets. Let \( x_{0} \) be an arbitrary object. It follows
    from elementary propositional logic that,
    \begin{equation*}
        x_{0} \in A \,\,{\land} \mleft( x_{0} \in B \land x_{0} \in C \mright)
        \iff \mleft( x_{0} \in A \land x_{0} \in B \mright) \land x_{0} \in C
    \end{equation*}
    By \cref{thm:intersection-in-terms-of-logical-and} the above statement is equivalent to,
    \begin{equation*}
        x_{0} \in A \intersect \mleft( B \intersect C \mright) \iff x_{0} \in \mleft( A \intersect B \mright) \intersect C
    \end{equation*}
    By Universal Generalization,
    \begin{equation*}
        \forall x \, \Big( x \in A \intersect \mleft( B \intersect C \mright)
        \iff x \in \mleft( A \intersect B \mright) \intersect C \Big)
    \end{equation*}
    By the \nameref{axm:axiom-of-extensionality}, \( A \intersect \mleft( B \intersect C \mright) = \mleft( A \intersect B \mright) \intersect C \).
    Since the sets \( A \), \( B \), and \( C \) are arbitrary, the theorem follows by Universal Generalization.
\end{proof}
\subsection{Difference}
\begin{thm}
    \label{set-identities:subtracing-empty-set-from-any-set-yields-the-set-itself}
    For any set \( A \), \( A - \emptyset = A \).
    \begin{equation*}
        \forall A \, \mleft( A - \emptyset = A \mright)
    \end{equation*}
\end{thm}

\begin{proof}
    Let \( A \) be an arbitrary set. Let \( x_{0} \) be an arbitrary object. By \cref{dfn:difference},
    \begin{equation*}
        x_{0} \in A - \emptyset \iff x_{0} \in A \land x_{0} \notin \emptyset
    \end{equation*}
    But \cref{dfn:empty-set} implies \( x_{0} \notin \emptyset \), thus,
    \begin{equation*}
        x_{0} \in A \land x_{0} \notin \emptyset \iff x_{0} \in A
    \end{equation*}
    Combining the two statements, we have,
    \begin{equation*}
        x_{0} \in A - \emptyset \iff x_{0} \in A
    \end{equation*}
    Since \( x_{0} \) is arbitrary, we conclude by Universal Generalization that,
    \begin{equation*}
        \forall x \, \mleft( x \in A - \emptyset \iff x \in A \mright)
    \end{equation*}
    By the \nameref{axm:axiom-of-extensionality}, it follows that \( A - \emptyset = A \).
\end{proof}

\begin{thm}
    \label{set-identities:subtracting-any-set-from-empty-set-is-still-empty}
    For any set \( A \), \( \emptyset - A = \emptyset \)
    \begin{equation*}
        \forall A \, \mleft( \emptyset - A = \emptyset \mright)
    \end{equation*}
\end{thm}

\begin{proof}
    Let \( A \) be an arbitrary set. Let \( x_{0} \) be an arbitrary object.
    \begin{equation*}
        \begin{aligned}
            x_{0} \notin \emptyset& \\
            x_{0} \notin \emptyset \lor x_{0} \in A& \\
            \lnot \, \mleft( x_{0} \in \emptyset \land x_{0} \notin A \mright)& \\
            x_{0} \notin \emptyset - A& \\
            \forall x \, \mleft( x \notin \emptyset - A \mright)&\\
            \emptyset - A = \emptyset&
        \end{aligned}
        \text{  }
        \begin{aligned}
            &\quad \text{\cref{dfn:empty-set}}\\
            &\quad \text{Disjunction Introduction}\\
            &\quad \text{De Morgan's Law}\\
            &\quad \text{\cref{dfn:difference}}\\
            &\quad \text{Universal Generalization}\\
            &\quad \text{\cref{dfn:empty-set}}
        \end{aligned}
    \end{equation*}
    Since \( A \) is arbitrary, the theorem follows from Universal Generalization.
\end{proof}

\begin{thm}[Nilpotence]
    \label{set-identities:nilpotence-of-relative-complements}
    For any set \( A \), \( A - A = \emptyset \)
    \begin{equation*}
        \forall A \, \mleft( A - A = \emptyset \mright)
    \end{equation*}
\end{thm}

\begin{proof}
    Let \( A \) be an arbitrary set. Let \( x_{0} \) be an arbitrary object.
    \begin{equation*}
        \begin{aligned}
            x_{0} \in A \lor x_{0} \notin A& \\
            \lnot \, \mleft( x_{0} \notin A \land x_{0} \in A \mright)& \\
            \lnot \, \mleft( x_{0} \in A \land x_{0} \notin A \mright)& \\
            x_{0} \notin A - A& \\
            \forall x \, \mleft( x \notin A - A \mright)& \\
            A - A = \emptyset&
        \end{aligned}
        \text{  }
        \begin{aligned}
            &\quad \text{Law of Excluded Middle}\\
            &\quad \text{De Morgan's Law}\\
            &\quad \text{\( \land \)-Commutativity}\\
            &\quad \text{\cref{dfn:difference}}\\
            &\quad \text{Universal Generalization}\\
            &\quad \text{\cref{dfn:empty-set}}
        \end{aligned}
    \end{equation*}
    Since \( A \) is arbitrary, the theorem follows from Universal Generalization.
\end{proof}

\begin{thm}
    \label{set-identities:subtract-B-minus-C-from-A}
    For any three sets \( A \), \( B \), and \( C \),
    \( A - \mleft( B - C \mright) = \mleft( A - B \mright) \union \mleft( A \intersect C \mright) \)
    \begin{equation*}
        \forall A,B,C \, \Big( A - \mleft( B - C \mright) = \mleft( A - B \mright) \union \mleft( A \intersect C \mright)  \Big)
    \end{equation*}
\end{thm}

\begin{proof}
    Let \( A \), \( B \), \( C \) be arbitrary sets. Let \( x_{0} \) be an arbitrary object.
    The goal is to use elementary propositional logic to re-express \( x_{0} \in A - \mleft( B - C \mright) \)
    in a logically equivalent form:
    \begin{align*}
        x_{0} \in A - \mleft( B - C \mright)
        &\iff
        x_{0} \in A \land x_{0} \notin \mleft( B - C \mright) \\
        &\iff
        x_{0} \in A \land \lnot \, \mleft( x_{0} \in B \land x_{0} \notin C \mright)\\
        &\iff
        x_{0} \in A \,\,{\land} \mleft( x_{0} \notin B \lor x_{0} \in C \mright)\\
        &\iff
        \mleft( x_{0} \in A \land x_{0} \notin B \mright) \lor
        \mleft( x_{0} \in A \land x_{0} \in C \mright)\\
        &\iff
        x_{0} \in \mleft( A - B \mright) \lor x_{0} \in \mleft( A \intersect C \mright) \\
        &\iff
        x_{0} \in \mleft( A - B \mright) \union \mleft( A \intersect C \mright)
    \end{align*}
    Since \( x_{0} \) is arbitrary, we conclude by Universal Generalization that,
    \begin{equation*}
        \forall x \, \Big( x \in A - \mleft( B - C \mright) \iff
        x \in \mleft( A - B \mright) \union \mleft( A \intersect C \mright) \Big)
    \end{equation*}
    By the \nameref{axm:axiom-of-extensionality}, \( A - \mleft( B - C \mright) =
    \mleft( A - B \mright) \union \mleft( A \intersect C \mright) \). Since \( A \), \( B \),
    and \( C \) are arbitrary, the theorem follows from Universal Generalization.
\end{proof}

\begin{thm}
    \label{set-identities:subtract-C-from-A-minus-B}
    For any three sets \( A \), \( B \), and \( C \), \( \mleft( A - B \mright) - C = A - \mleft( B \union C \mright) \)
    \begin{equation*}
        \forall A,B,C , \Big( \mleft( A - B \mright) - C = A - \mleft( B \union C \mright) \Big)
    \end{equation*}
\end{thm}

\begin{proof}
    Let \( A \), \( B \), and \( C \) be arbitrary sets. Let \( x_{0} \) be an arbitrary object. Again, we use
    the laws of propositional logic, \cref{dfn:difference,thm:union-in-terms-of-logical-or} to express
    \( x_{0} \in \mleft( A - B \mright) - C \) in another logically equivalent form:
    \begin{align*}
        x_{0} \in \mleft( A - B \mright) - C
        &\iff
        x_{0} \in \mleft( A - B \mright) \land x_{0} \notin C \\
        &\iff
        \mleft( x_{0} \in A \land x_{0} \notin B \mright)
        \land
        x_{0} \notin C\\
        &\iff
        x_{0} \in A \,\,{\land} \mleft( x_{0} \notin B \land x_{0} \notin C \mright)\\
        &\iff
        x_{0} \in A \land \lnot \, \mleft( x_{0} \in B \lor x_{0} \in C \mright)\\
        &\iff
        x_{0} \in A \land x_{0} \notin B \union C\\
        &\iff
        x_{0} \in A - \mleft( B \union C \mright)
    \end{align*}
    Since \( x_{0} \) is arbitrary, we conclude by Universal Generalization that,
    \begin{equation*}
        \forall x \, \Big( x \in \mleft( A - B \mright) - C \iff
        x \in A - \mleft( B \union C \mright) \Big)
    \end{equation*}
    By the \nameref{axm:axiom-of-extensionality}, \( \mleft( A - B \mright) - C = A - \mleft( B \union  C \mright)\).
    Since \( A \), \( B \), and \( C \) are arbitrary, the theorem follows from Universal Generalization.
\end{proof}

\begin{thm}[Nilpotence]
    \label{set-identities:symmetric-difference-of-any-set-with-itself-is-empty}
    For any set \( A \), \( A \symdiff A = \emptyset \).
\end{thm}

\begin{proof}
    Let \( A \) be an arbitrary set. By \cref{dfn:symmetric-difference},
    \begin{equation*}
        \forall x \, \Big( x \in A \symdiff A \iff
        \mleft( x \in A \lor x \in A \mright) \land
        \lnot \, \mleft( x \in A \land x \in A \mright)\Big)
    \end{equation*}
    Using the laws of propositional logic, we can transform the statement into
    another logically equivalent form:
    \begin{equation*}
        \forall x \, \Big( x \in A \symdiff A \iff
        x \in A \land x \notin A \Big)
    \end{equation*}
    But by the Law of Non-Contradiction, the statement on the right-hand side of the biconditional is always false.
    \begin{equation*}
        \therefore \, \forall x \, \mleft( x \notin A \symdiff A \mright)
    \end{equation*}
    By \cref{dfn:empty-set}, \( A \symdiff A = \emptyset \).
\end{proof}

\begin{thm}
    \label{set-identities:symmetric-difference-in-terms-of-union-and-intersection}
    For any two sets \( A \) and \( B \), \( A \symdiff B = \mleft( A \union B \mright) - \mleft( A \intersect B \mright) \)
    \begin{equation*}
        \forall A,B \, \Big( A \symdiff B = \mleft( A \union B \mright) - \mleft( A \intersect B \mright) \Big)
    \end{equation*}
\end{thm}

\begin{proof}
    Let \( A \) and \( B \) be arbitrary sets. Let \( x_{0} \) be an arbitrary object. We use the laws of logic,
    \cref{dfn:symmetric-difference,thm:union-in-terms-of-logical-or,thm:intersection-in-terms-of-logical-and}
    to re-express \( x_{0} \in A \symdiff B \) in another logically equivalent form:
    \begin{align*}
        x_{0} \in A \symdiff B
        &\iff
        \mleft( x_{0} \in A \lor x_{0} \in B \mright)
        \land
        \lnot \, \mleft( x_{0} \in A \land x_{0} \in B \mright) \\
        &\iff
        x_{0} \in A \union B \land x_{0} \notin A \intersect B\\
        &\iff
        x_{0} \in \big( \mleft( A \union B \mright) - \mleft( A \intersect B \mright) \big)
    \end{align*}
    Since \( x_{0} \) is arbitrary, we conclude by Universal Generalization that,
    \begin{equation*}
        \forall x \, \Big( x \in A \symdiff B \iff
        x \in \big( \mleft( A \union B \mright) - \mleft( A \intersect B \mright) \big)\Big)
    \end{equation*}
    By the \nameref{axm:axiom-of-extensionality}, \( A \symdiff B = \mleft( A \union B \mright) - \mleft( A \intersect B \mright) \).
    Since \( A \), \( B \), and \( C \) are arbitrary, the theorem follows by Universal Generalization.
\end{proof}

\begin{thm}
    \label{set-identities:symmetric-difference-as-the-union-of-differences}
    For any two sets \( A \) and \( B \), \( A \symdiff B = \mleft( A - B \mright) \union \mleft( B - A \mright) \)
    \begin{equation*}
        \forall A,B \, \Big( A \symdiff B = \mleft( A - B \mright) \union \mleft( B - A \mright) \Big)
    \end{equation*}
\end{thm}

\begin{proof}
    Let \( A \) and \( B \) be arbitrary sets. Let \( x_{0} \) be an arbitrary object.
    Again, we use the laws of propositional logic and known definitions to re-express
    \( x_{0} \in A \symdiff B \) in an alternative form:
    \begin{align*}
        x_{0} \in A \symdiff B
        &\iff
        \mleft( x_{0} \in A \lor x_{0} \in B \mright)
        \land
        \lnot \, \mleft( x_{0} \in A \land x_{0} \in B \mright)\\
        &\iff
        \mleft( x_{0} \in A \lor x_{0} \in B \mright)
        \land
        \mleft( x_{0} \notin A \lor x_{0} \notin B \mright)\\
        &\iff
        \big( \mleft( x_{0} \in A \lor x_{0} \in B \mright) \land x_{0} \notin A \big)
        \lor
        \big( \mleft( x_{0} \in A \lor x_{0} \in B \mright) \land x_{0} \notin B \big)\\
        &\iff
        \mleft( x_{0} \in B \land x_{0} \notin A \mright)
        \lor
        \mleft( x_{0} \in A \land x_{0} \notin B \mright)\\
        &\iff
        x_{0} \in \mleft( A - B \mright) \lor x_{0} \in \mleft( B - A \mright)\\
        &\iff
        x_{0} \in \big( \mleft( A - B \mright) \union \mleft( B - A \mright) \big)
    \end{align*}
    Since \( x_{0} \) is arbitrary, we have,
    \begin{equation*}
        \forall x \, \Big( x \in A \symdiff B \iff x \in \big( \mleft( A - B \mright) \union \mleft( B - A \mright) \big) \Big)
    \end{equation*}
    By the \nameref{axm:axiom-of-extensionality}, \( A \symdiff B = \mleft( A - B \mright) \union \mleft( B - A \mright) \).
    Since \( A \) and \( B \) are arbitrary, the theorem follows from Universal Generalization.
\end{proof}

\begin{thm}[Commutativity]
    \label{set-identities:commutativity-of-symmetric-difference}
    For any two sets \( A \) and \( B \), \( A \symdiff B = B \symdiff A \).
    \begin{equation*}
        \forall A,B \, \mleft( A \symdiff B = B \symdiff A \mright)
    \end{equation*}
\end{thm}

\begin{proof}
    Let \( A \) and \( B \) be arbitrary sets. By \cref{set-identities:symmetric-difference-as-the-union-of-differences,set-identities:commutativity-of-union},
    \begin{align*}
        A \symdiff B
        &=
        \mleft( A - B \mright) \union \mleft( B - A \mright)\\
        &=
        \mleft( B - A \mright) \union \mleft( A - B \mright)\\
        &=
        B \symdiff A
    \end{align*}
    Since \( A \) and \( B \) are arbitrary, the theorem follows from Universal Generalization.
\end{proof}

\begin{thm}[Associativity]
    \label{set-identities:associativity-of-symmetric-difference}
    For any three sets \( A \), \( B \), and \( C \),
    \begin{equation*}
        A \symdiff \mleft( B \symdiff C \mright) =
        \mleft( A \symdiff B \mright) \symdiff C
    \end{equation*}
\end{thm}

\begin{proof}
    Let \( A \), \( B \), and \( C \) be arbitrary sets. Let \( x_{0} \) be an arbitrary object.
    Furthermore, suppose that \( x_{0} \in A \symdiff \mleft( B \symdiff C \mright) \).
    Then by \cref{dfn:symmetric-difference}, it follows that,
    \begin{equation}
        \label{eq:lhs-predicate-form}
        \big( x_{0} \in A \land x_{0} \notin \mleft( B \symdiff C \mright) \big)
        \lor
        \big( x_{0} \notin A \land x_{0} \in \mleft( B \symdiff C \mright) \big)
    \end{equation}
    Suppose that \( x_{0} \in A \land x_{0} \notin \mleft( B \symdiff C \mright) \) is true.
    Using \cref{dfn:symmetric-difference} once again, we infer,
    \begin{equation*}
        \mleft( x_{0} \in A \land x_{0} \in B \land x_{0} \in C \mright)
        \lor
        \mleft( x_{0} \in A \land x_{0} \notin B \land x_{0} \notin C \mright)
    \end{equation*}
    If \( x_{0} \in A \,\,{\land} \mleft( x_{0} \in B \land x_{0} \in C \mright)\),
    then clearly \( x_{0} \in \mleft( A \symdiff B \mright) \symdiff C \) because
    \( x_{0} \in C \) but \( x_{0} \notin \mleft( A \symdiff B \mright) \). Similarly, if
    \( x_{0} \in A \,\,{\land} \mleft( x_{0} \notin B \land x_{0} \notin C \mright) \),
    then \( x_{0} \in \mleft( A \symdiff B \mright) \symdiff C \) also follows because
    \( x_{0} \in \mleft( A \symdiff B \mright) \) but \( x_{0} \notin C \). Since both disjuncts
    imply \( x_{0} \in \mleft( A \symdiff B \mright) \symdiff C \), we apply Disjunction
    Elimination to infer,
    \begin{equation}
        \label{eq:left-disjunct-forward-implication}
        x_{0} \in A \land x_{0} \notin \mleft( B \symdiff C \mright)
        \implies
        x_{0} \in \mleft( A \symdiff B \mright) \symdiff C
    \end{equation}

    \vspace{0.1cm}
    \noindent Similarly, suppose that \( x_{0} \notin A \land x_{0} \in \mleft( B \symdiff C \mright) \) is true.
    Applying \cref{dfn:symmetric-difference}, we infer,
    \begin{equation*}
        \mleft( x_{0} \notin A \land x_{0} \in B \land x_{0} \notin C \mright)
        \lor
        \mleft( x_{0} \notin A \land x_{0} \notin B \land x_{0} \in C \mright)
    \end{equation*}
    If \( x_{0} \notin A \land x_{0} \in B \land x_{0} \notin C \),
    then clearly \( x_{0} \in \mleft( A \symdiff B \mright) \symdiff C \) because
    \( x_{0} \in \mleft( A \symdiff B \mright) \) but \( x_{0} \notin C \).
    If \( x_{0} \notin A \land x_{0} \notin B \land x_{0} \in C \),
    then \( x_{0} \in \mleft( A \symdiff B \mright) \symdiff C \) is also true because
    \( x_{0} \notin \mleft( A \symdiff B \mright) \) but \( x_{0} \in C \). Since both disjuncts
    imply \( x_{0} \in \mleft( A \symdiff B \mright) \symdiff C \), Disjunction Elmination
    yields,
    \begin{equation}
        \label{eq:right-disjunct-forward-implication}
        x_{0} \notin A \land x_{0} \in \mleft( B \symdiff C \mright)
        \implies
        x_{0} \in \mleft( A \symdiff B \mright) \symdiff C
    \end{equation}
    Applying Disjunction Elimination to \eqref{eq:lhs-predicate-form}, \eqref{eq:left-disjunct-forward-implication},
    and \eqref{eq:right-disjunct-forward-implication}, we infer,
    \begin{equation}
        \label{eq:symmetric-difference-associativity-forward-implication}
        x_{0} \in A \symdiff \mleft( B \symdiff C \mright)
        \implies
        x_{0} \in \mleft( A \symdiff B \mright) \symdiff C
    \end{equation}

    \vspace{0.1cm}
    \noindent To prove the converse, we suppose that \( x_{0} \in \mleft( A \symdiff B \mright) \symdiff C \).
    Then by \cref{dfn:symmetric-difference}, it follows that,
    \begin{equation}
        \label{eq:rhs-predicate-form}
        \big( x_{0} \in \mleft( A \symdiff B \mright) \land x_{0} \notin C \big)
        \lor
        \big( x_{0} \notin \mleft( A \symdiff B \mright) \land x_{0} \in C \big)
    \end{equation}
    Suppose that \( x_{0} \in \mleft( A \symdiff B \mright) \land x_{0} \notin C \). It follows from
    \cref{dfn:symmetric-difference} that,
    \begin{equation*}
        \mleft( x_{0} \in A \land x_{0} \notin B \land x_{0} \notin C \mright)
        \lor
        \mleft( x_{0} \notin A \land x_{0} \in B \land x_{0} \notin C \mright)
    \end{equation*}
    If \( \mleft( x_{0} \in A \land x_{0} \notin B \land x_{0} \notin C \mright) \),
    then \( x_{0} \in A \symdiff \mleft( B \symdiff C \mright) \) because \( x_{0} \in A \) and
    \( x_{0} \notin \mleft( B \symdiff C \mright) \).
    Similarly, if \( \mleft( x_{0} \notin A \land x_{0} \in B \land x_{0} \notin C \mright) \),
    then \( x_{0} \in A \symdiff \mleft( B \symdiff C \mright) \) is also true because
    \( x_{0} \notin A \) and \( x_{0} \in \mleft( B \symdiff C \mright) \). Since both disjuncts imply
    \( x_{0} \in A \symdiff \mleft( B \symdiff C \mright) \), we conclude by Disjunction Elimination that,
    \begin{equation}
        \label{eq:left-disjunct-backward-implication}
        x_{0} \in \mleft( A \symdiff B \mright) \land  x_{0} \notin C
        \implies
        x_{0} \in A \symdiff \mleft( B \symdiff C \mright)
    \end{equation}
    Next, suppose that \( x_{0} \notin \mleft( A \symdiff B \mright) \land \, x_{0} \in C \) is true. Once again
    applying \cref{dfn:symmetric-difference}, we have,
    \begin{equation*}
        \mleft( x_{0} \notin A \land x_{0} \notin B \land x_{0} \in C \mright)
        \lor
        \mleft( x_{0} \in A \land x_{0} \in B \land x_{0} \in C \mright)
    \end{equation*}
    If \( \mleft( x_{0} \notin A \land x_{0} \notin B \land x_{0} \in C \mright) \),
    then \( x_{0} \in A \symdiff \mleft( B \symdiff C \mright) \) must hold because
    \( x_{0} \notin A \) and \( x_{0} \in \mleft( B \symdiff C \mright) \).
    If \( \mleft( x_{0} \in A \land x_{0} \in B \land x_{0} \in C \mright) \), then
    \( x_{0} \in A \symdiff \mleft( B \symdiff C \mright) \) is also true because \( x_{0} \in A \)
    and \( x_{0} \notin \mleft( B \symdiff C \mright) \). Since both disjuncts imply \( x_{0} \in
    A \symdiff \mleft( B \symdiff C \mright)\), we conclude by Disjunction Elimination that,
    \begin{equation}
        \label{eq:right-disjunct-backward-implication}
        x_{0} \notin \mleft( A \symdiff B \mright) \land \,\, x_{0} \in C
        \implies
        x_{0} \in A \mleft( B \symdiff C \mright)
    \end{equation}
    Applying Disjunction Elimination to \eqref{eq:rhs-predicate-form}, \eqref{eq:left-disjunct-backward-implication},
    and \eqref{eq:right-disjunct-backward-implication} yields,
    \begin{equation}
        \label{eq:symmetric-difference-associativity-backward-implication}
        x_{0} \in \mleft( A \symdiff B \mright) \symdiff C
        \implies
        x_{0} \in A \symdiff \mleft( B \symdiff C \mright)
    \end{equation}
    Together, \eqref{eq:symmetric-difference-associativity-forward-implication} and \eqref{eq:symmetric-difference-associativity-backward-implication}
    establish the biconditional:
    \begin{equation*}
        x_{0} \in A \symdiff \mleft( B \symdiff C \mright)
        \iff
        x_{0} \in \mleft( A \symdiff B \mright) \symdiff C
    \end{equation*}
    Since \( x_{0} \) is arbitrary, Universal Generalization yields,
    \begin{equation*}
        \forall x \, \Big( x \in A \symdiff \mleft( B \symdiff C \mright)
        \iff
        x \in \mleft( A \symdiff B \mright) \symdiff C
        \Big)
    \end{equation*}
    Since \( A \), \( B \), and \( C \) are arbitrary, the theorem follows from
    Universal Generalization.
\end{proof}
\subsection{Mixed Operations}
\begin{thm}[Generalized De Morgan Laws]
    \label{set-identities:generalized-de-morgan-laws}
    For any set \( A \) and any non-empty system of sets \( S \),
    \begin{gather*}
        A - \Union S =
        \Intersect \setbuilderbig{x}{\exists z \, \mleft( z \in S \land x = A - z \mright)}\\
        A - \Intersect S =
        \Union \setbuilderbig{x}{\exists z \, \mleft( z \in S \land x = A - z \mright)}
    \end{gather*}
\end{thm}

\begin{proof}
    Let \( A \) be an arbitrary set. Let \( S \) be an arbitrary system of sets such that \( S \neq \emptyset \).
    Let the set \( T \) be defined as,
    \begin{equation}
        \label{eq:T-in-set-builder-notation}
        T = \setbuilderbig{x}{\exists z \, \mleft( z \in S \land x = A - z \mright)}
    \end{equation}
    First, we must prove that the set \( T \) does indeed exist. Let \( x_{0} \) be an arbitrary object such that
    \( \exists z \, \mleft( z \in S \land x_{0} = A - z \mright) \). In other words, \( x_{0} = A - z_{0} \) for
    some \( z_{0} \in S \). By \cref{dfn:difference}, it follows that,
    \begin{equation*}
        \forall x \, \mleft( x \in x_{0} \iff x \in A \land x \notin z_{0} \mright)
    \end{equation*}
    But \( \forall x \, \mleft( x \in A \land x \notin z_{0} \implies x \in A \mright) \) is a tautology.
    \begin{equation*}
        \therefore \, \forall x \, \mleft( x \in x_{0} \implies x \in A \mright)
    \end{equation*}
    In other words, \( x_{0} \subseteq A \). Equivalently, \( x_{0} \in \mathcal{P}(A) \).
    \begin{equation*}
        \therefore \, \forall x \, \mleft( \exists z \, \mleft( z \in S \land x = A - z \mright) \implies x \in \mathcal{P}(A) \mright)
    \end{equation*}
    The existence of the set \( T \) is thus guaranteed by \cref{thm:sufficient-condition-for-unrestricted-comprehension}.

    \vspace{0.5cm}
    \noindent
    Using \cref{dfn:intersection-of-sets}, we can write out the definition for \( \Intersect T \).
    \begin{equation}
        \label{eq:statement-for-intersection-of-T}
        \forall x \, \mleft( x \in \Intersect T
        \iff
        \forall y \, \mleft( y \in T \implies x \in y \mright) \mright)
    \end{equation}
    Applying \eqref{eq:T-in-set-builder-notation} to \eqref{eq:statement-for-intersection-of-T} yields,
    \begin{equation}
        \label{eq:statement-for-intersection-of-T-substituted}
        \forall x \, \mleft( x \in \Intersect T
        \iff
        \forall y \, \mleft( \exists z \, \mleft( z \in S \land y = A - z \mright) \implies x \in y \mright)\mright)
    \end{equation}
    Let \( x_{0} \) be an arbitrary object. Applying Universal Instantiation to
    \eqref{eq:statement-for-intersection-of-T-substituted} yields,
    \begin{equation}
        \label{eq:instantiated-statement-for-intersection-of-T-substituted}
        x_{0} \in \Intersect T
        \iff
        \forall y \, \mleft( \exists z \, \mleft( z \in S \land y = A - z \mright) \implies x_{0} \in y \mright)
    \end{equation}
    Using elementary logic, we can re-express the predicate on the right-hand side of the biconditional
    in various equivalent forms:
    \begin{equation*}
        \begin{aligned}
            &\forall y \, \mleft( \exists z \, \mleft( z \in S \land
            y = A - z \mright) \implies x_{0} \in y \mright) \\
            &\forall y \, \mleft( \lnot \, \exists z \, \mleft( z \in S \land y = A - z \mright) \lor x_{0} \in y \mright) \\
            &\forall y \mleft( \forall z \, \mleft( z \notin S \lor y \neq A - z \mright) \lor x_{0} \in y \mright) \\
            &\forall y \, \forall z \, \mleft( z \notin S \lor y \neq A - z \lor x_{0} \in y \mright) \\
            &\forall z \, \mleft( z \notin S \lor \forall y \, \mleft( y \neq A - z \lor x_{0} \in y \mright) \mright)\\
            &\forall z \, \mleft( z \notin S \lor \forall y \, \mleft( y = A - z \implies x_{0} \in y \mright) \mright)\\
            &\forall z \, \mleft( z \notin S \lor x_{0} \in A - z \mright)\\
            &\forall z \, \mleft( z \notin S \,\,{\lor} \mleft( x_{0} \in A \land x_{0} \notin z \mright) \mright)\\
            &\forall z \, \mleft( z \notin S \lor x_{0} \in A \land z \notin S \lor x_{0} \notin z \mright)\\
            &\forall z \, \mleft( z \notin S \lor x_{0} \in A \mright)
            \land
            \forall z \, \mleft( z \notin S \lor x_{0} \notin z \mright)\\
            &x_{0} \in A \lor \forall z \, \mleft( z \notin S \mright) \land
            \forall z \, \mleft( z \notin S \lor x_{0} \notin z \mright)\\
            &x_{0} \in A \lor \forall z \, \mleft( z \notin S \mright) \land
            \forall z \, \mleft( z \in S \implies x_{0} \notin z \mright)\\
        \end{aligned}
        \text{  }
        \begin{aligned}
            &\quad \text{by \eqref{eq:statement-for-intersection-of-T-substituted}} \\
            &\quad \text{Definition of Conditional}\\
            &\quad \text{De Morgan's Law}\\
            &\quad \text{\cref{logic:grouping-disjuncts-into-quantified-statement-universal}}\\
            &\quad \text{\cref{logic:grouping-disjuncts-into-quantified-statement-universal}}\\
            &\quad \text{Definition of Conditional}\\
            &\quad \text{\cref{logic:using-universal-quantifier-to-express-the-fact-that-an-object-satisfies-a-predicate}}\\
            &\quad \text{\cref{dfn:difference}}\\
            &\quad \text{\( \lor \)-Distributivity}\\
            &\quad \text{Logical Equivalence}\\
            &\quad \text{\cref{logic:grouping-disjuncts-into-quantified-statement-universal}}\\
            &\quad \text{Definition of Conditional}\\
        \end{aligned}
    \end{equation*}
    Thus, we have,
    \begin{equation*}
        x_{0} \in \Intersect T
        \iff
        x_{0} \in A \lor  \forall z \, \mleft( z \notin S \mright) {\land}\,\,  \forall z \, \mleft( z \in S \implies x_{0} \notin z \mright)
    \end{equation*}
    Since \( S \neq \emptyset \), \( \forall z \, \mleft( z \notin S \mright) \) must be false. Thus, the formula reduces to,
    \begin{equation}
        \label{eq:final-statement-for-intersect-T}
        x_{0} \in \Intersect T
        \iff
        x_{0} \in A \land
        \forall z \, \mleft( z \in S \implies x_{0} \notin z \mright)
    \end{equation}
    By \cref{dfn:difference,dfn:union-of-sets},
    \begin{align}
        \label{eq:final-statement-for-A-minus-union-S}
        x_{0} \in A - \Union S
        &\iff
        x_{0} \in A \land x_{0} \notin \Union S \notag \\
        &\iff
        x_{0} \in A \land \lnot \, \exists z \, \mleft( z \in S \land x_{0} \in z \mright) \notag \\
        &\iff
        x_{0} \in A \land \forall z \, \mleft( z \notin S \lor x_{0} \notin z \mright)\notag \\
        &\iff
        x_{0} \in A \land \forall z \, \mleft( z \in S \implies x_{0} \notin z \mright)
    \end{align}
    Comparing \eqref{eq:final-statement-for-intersect-T} and \eqref{eq:final-statement-for-A-minus-union-S}, we conclude that,
    \begin{equation*}
        x_{0} \in A - \Union S \iff x_{0} \in \Intersect T
    \end{equation*}
    Since \( x_{0} \) is arbitrary, Universal Generalization yields,
    \begin{equation*}
        \forall x \, \mleft( x \in A - \Union S \iff x \in \Intersect T \mright)
    \end{equation*}
    By the \nameref{axm:axiom-of-extensionality}, it follows that,
    \begin{equation*}
        A - \Union S = \Intersect T
    \end{equation*}
    To prove the second identity, we start with the definition of the union of sets
    given by \cref{dfn:union-of-sets}.
    \begin{equation}
        \label{eq:statement-for-union-of-T}
        \forall x \, \mleft( x \in \Union T
        \iff
        \exists y \, \mleft( y \in T \land x \in y \mright)\mright)
    \end{equation}
    Applying \eqref{eq:T-in-set-builder-notation} to \eqref{eq:statement-for-union-of-T},
    \begin{equation}
        \label{eq:statement-for-union-of-T-substituted}
        \forall x \, \mleft( x \in \Union T
        \iff
        \exists y \, \mleft( \exists z \, \mleft( z \in S \land y = A - z \mright)
        \land
        x \in y \mright)\mright)
    \end{equation}
    Let \( x_{0} \) be an arbitrary object.
    From \eqref{eq:statement-for-union-of-T-substituted}, we infer by Universal Instantiation that,
    \begin{equation}
        \label{eq:instantiated-statement-for-union-of-T-substituted}
        x_{0} \in \Union T
        \iff
        \exists y \, \mleft( \exists z \, \mleft( z \in S \land y = A - z \mright)
        \land
        x_{0} \in y \mright)
    \end{equation}
    Once again, we apply the laws of elementary logic to transform the predicate on the right-hand side of
    the biconditional into various logically equivalent forms:
    \begin{equation*}
        \begin{aligned}
            &\exists y \, \mleft( \exists z \, \mleft( z \in S \land y = A - z \mright) \land x_{0} \in y \mright)\\
            &\exists y \, \exists z \, \mleft( z \in S \land y = A - z \land x_{0} \in y \mright)\\
            &\exists z \, \mleft( z \in S \land \exists y \, \mleft( y = A - z \land x_{0} \in y \mright) \mright)\\
            &\exists z \, \mleft( z \in S \land x_{0} \in A - z \mright)\\
            &\exists z \, \mleft( z \in S \land \mleft( x_{0} \in A \land x_{0} \notin z \mright) \mright)\\
            &\exists z \, \mleft( x_{0} \in A \land \mleft( z \in S \land x_{0} \notin z \mright) \mright)\\
            &x_{0} \in A \land \exists z \, \mleft( z \in S \land x_{0} \notin z \mright)\\
        \end{aligned}
        \text{  }
        \begin{aligned}
            &\quad \text{from \eqref{eq:instantiated-statement-for-union-of-T-substituted}}\\
            &\quad \text{\cref{logic:grouping-conjuncts-into-quantified-statement-existential}}\\
            &\quad \text{\cref{logic:grouping-conjuncts-into-quantified-statement-existential}}\\
            &\quad \text{\cref{logic:using-existential-quantifier-to-express-the-fact-that-an-object-satisfies-a-predicate}}\\
            &\quad \text{\cref{dfn:difference}}\\
            &\quad \text{\( \land \)-Associativity}\\
            &\quad \text{\cref{logic:grouping-conjuncts-into-quantified-statement-existential}}\\
        \end{aligned}
    \end{equation*}
    Thus, we have,
    \begin{equation}
        \label{eq:final-statement-for-union-T}
        x_{0} \in \Union T
        \iff
        x_{0} \in A \land \exists z \, \mleft( z \in S \land x_{0} \notin z \mright)
    \end{equation}
    By \cref{dfn:intersection-of-sets,dfn:difference},
    \begin{align}
        \label{eq:final-statement-for-A-minus-intersect-S}
        x_{0} \in A - \Intersect S
        &\iff
        x_{0} \in A \land x_{0} \notin \Intersect S \notag \\
        &\iff
        x_{0} \in A \land \lnot \, \forall z \, \mleft( z \in S \implies x_{0} \in z \mright) \notag \\
        &\iff
        x_{0} \in A \land \lnot \, \forall z \, \mleft( z \notin S \lor x_{0} \in z \mright) \notag \\
        &\iff
        x_{0} \in A \land \exists z \, \mleft( z \in S \land x_{0} \notin z \mright)
    \end{align}
    From \eqref{eq:final-statement-for-union-T} and \eqref{eq:final-statement-for-A-minus-intersect-S}, we conclude that,
    \begin{equation*}
        x_{0} \in A - \Intersect S \iff x_{0} \in \Union T
    \end{equation*}
    And since \( x_{0} \) is arbitrary, Universal Generalization yields,
    \begin{equation*}
        \forall x \, \mleft( x \in A - \Intersect S \iff x \in \Union T \mright)
    \end{equation*}
    By the \nameref{axm:axiom-of-extensionality},
    \begin{equation*}
        A - \Intersect S = \Union T \qedhere
    \end{equation*}
\end{proof}

\begin{thm}[De Morgan Laws]
    \label{set-identities:de-morgan-laws}
    For any three sets \( A \), \( B \), and \( C \),
    \begin{gather*}
        A - \mleft( B \union C \mright) = \mleft( A - B \mright) \intersect \mleft( A - C \mright) \\
        A - \mleft( B \intersect C \mright) = \mleft( A - B \mright) \union \mleft( A - C \mright)
    \end{gather*}
\end{thm}

\begin{proof}
    This is merely a special case of \cref{set-identities:generalized-de-morgan-laws} with
    \( S = \mleft\{ B, C \mright\} \).
    Thus, by \cref{set-identities:generalized-de-morgan-laws,dfn:union-of-sets,dfn:unordered-pair,dfn:intersection-of-sets},
    \begin{align*}
        A - \mleft( B \union C \mright)
        &=
        A - \Union \Big\{ B, C \Big\} \\
        &=
        \Intersect \setbuilderbig{x \in \powerset{A}}{\exists z \, \mleft( z \in \mleft\{ A, B\mright\} \land x = A - z \mright)}\\
        &=
        \Intersect \Big\{ A - B, A - C \Big\} \\
        &=
        \mleft( A - B \mright) \intersect \mleft( A - C \mright)
    \end{align*}
    Similarly,
    \begin{align*}
        A - \mleft( B \intersect C \mright)
        &=
        A - \Intersect \Big\{ B, C \Big\} \\
        &=
        \Union \setbuilderbig{x \in \powerset{A}}{\exists z \, \mleft( z \in \mleft\{ B,C \mright\} \land x = A - z \mright)}\\
        &=
        \Union \Big\{ A - B, A - C \Big\} \\
        &=
        \mleft( A - B \mright) \union \mleft( A - C \mright) \qedhere
    \end{align*}
\end{proof}

\begin{thm}[Generalized Distributive Laws]
    \label{set-identities:generalized-distributive-laws}
    For any set \( A \) and any non-empty system of sets \( S \),
    \begin{gather*}
        A \union \Intersect S =
        \Intersect \setbuilderbig{x}{\exists z \, \mleft( z \in S \land x = A \union z \mright)}\\
        A \intersect \Union S =
        \Union \setbuilderbig{x}{\exists z \, \mleft( z \in S \land x = A \intersect z \mright)}
    \end{gather*}
\end{thm}

\begin{proof}
    Let \( A \) be a set and \( S \) be a non-empty system of sets. Let the sets \( T_{1} \) and \( T_{2} \)
    be defined as follows.
    \begin{gather}
        T_{1} =
        \setbuilderbig{x}{\exists z \, \mleft( z \in S \land x = A \union z \mright)} \label{eq:define-T1} \\
        T_{2} =
        \setbuilderbig{x}{\exists z \, \mleft( z \in S \land x = A \intersect z \mright)} \label{eq:define-T2}
    \end{gather}
    We begin by proving that \( T_{1} \) and \( T_{2} \) exist.

    \vspace{0.5cm}
    \noindent
    Let \( x_{0} \) be an arbitrary object such that
    \( \exists z \, \mleft( z \in S \land x_{0} = A \union z \mright) \) holds.
    In other words, \( x_{0} = A \union z_{0} \) for some \( z_{0} \in S \).
    Now consider the set \( A \union \Union S \), which exists by the \nameref{axm:axiom-of-union}.
    Using \cref{dfn:union-of-sets,thm:union-in-terms-of-logical-or}, we can express the definition
    of \( A \union \Union S \) like so:
    \begin{equation}
        \label{eq:define-A-union-the-union-of-S}
        \forall x \, \mleft( x \in A \union \Union S
        \iff
        x \in A \lor
        \exists y \, \mleft( y \in S \land x \in y \mright)\mright)
    \end{equation}
    Since \( x_{0} = A \union z_{0} \), it follows from \cref{thm:union-in-terms-of-logical-or}
    that,
    \begin{equation}
        \label{eq:predicate-for-A-union-z0}
        \forall x \, \mleft( x \in x_{0}
        \iff
        x \in A \lor x \in  z_{0} \mright)
    \end{equation}
    Let \( \alpha \) be an arbitrary object such that \( \alpha \in x_{0} \). By \eqref{eq:predicate-for-A-union-z0},
    \begin{equation*}
        \alpha \in A \lor \alpha \in z_{0}
    \end{equation*}
    If \( \alpha \in A \), then \( \alpha \in A \union \Union S \) by \eqref{eq:define-A-union-the-union-of-S}. If
    \( \alpha \in z_{0} \), then \( \exists y \, \mleft( y \in S \land \alpha \in y \mright) \) must be true
    because \( z_{0} \in S \). Thus, \( \alpha \in A \union \Union S \) also follows by \eqref{eq:define-A-union-the-union-of-S}.
    By Disjunction Elimination, we conclude that,
    \begin{equation*}
        \alpha \in x_{0} \implies \alpha \in A \union \Union S
    \end{equation*}
    Since \( \alpha \) is arbitrary, we can generalize,
    \begin{equation*}
        \forall x \, \mleft( x \in x_{0} \implies A \union \Union S \mright)
    \end{equation*}
    In other words, we have proven that \( x_{0} \subseteq A \union \Union S \) and thus \( x_{0} \in \mathcal{P}(A \union \Union S) \).
    This means that,
    \begin{equation*}
        \exists z \, \mleft( z \in S \land x_{0} = A \union z \mright) \implies x_{0} \in \mathcal{P}(A \union \Union S)
    \end{equation*}
    Generalizing once again,
    \begin{equation*}
        \forall x \, \mleft( \exists z \, \mleft( z \in S \land x = A \union z \mright) \implies x \in \mathcal{P}(A \union \Union S) \mright)
    \end{equation*}
    Therefore, \( T_{1} \) exists by \cref{thm:sufficient-condition-for-unrestricted-comprehension}.

    \vspace{0.5cm}
    \noindent
    Let \( x_{0} \) be an arbitrary object such that \( \exists z \, \mleft( z \in S \land x_{0} = A \intersect z \mright) \) holds.
    In other words, \( x_{0} = A \intersect z_{0} \) for some \( z_{0} \in S \).
    Clearly, \( x_{0} \subseteq A \), so \( x_{0} \in \mathcal{P}(A) \).
    By the Deduction Theorem, we conclude that,
    \begin{equation*}
        \exists z \, \mleft( z \in S \land x_{0} = A \intersect z \mright)
        \implies
        x_{0} \in \mathcal{P}(A)
    \end{equation*}
    By Univeral Generalization,
    \begin{equation*}
        \forall x \, \mleft( \exists z \, \mleft( z \in S \land x = A \intersect z \mright) \implies x \in \mathcal{P}(A) \mright)
    \end{equation*}
    The existence of \( T_{2} \) is guaranteed by \cref{thm:sufficient-condition-for-unrestricted-comprehension}.

    \vspace{0.5cm}
    \noindent
    Having established the existences of \( T_{1} \) and \( T_{2} \), we are now ready to prove that,
    \begin{gather*}
        A \union \Intersect S =
        \Intersect T_{1} \\
        A \intersect \Union S =
        \Union T_{2}
    \end{gather*}
    We do this by taking the definition of the set on the right-hand side of the equality and then
    transforming it (using the laws of elementary logic) into the definition of the set on the
    left-hand side.
    E.g: for the first equality, we apply \cref{dfn:intersection-of-sets} to the right-hand side to obtain,
    \begin{equation}
        \label{eq:statement-for-the-intersection-of-T1}
        \forall x \, \mleft( x \in \Intersect T_{1}
        \iff
        \forall y \, \mleft( y \in T_{1} \implies x \in y \mright)\mright)
    \end{equation}
    Next, we apply \eqref{eq:define-T1} to \eqref{eq:statement-for-the-intersection-of-T1},
    \begin{equation}
        \label{eq:statement-for-the-intersection-of-T1-substituted}
        \forall x \, \mleft( x \in \Intersect T_{1}
        \iff
        \forall y \, \mleft( \exists z \, \mleft( z \in S \land y = A \union z \mright) \implies x \in y \mright)\mright)
    \end{equation}
    Let \( x_{0} \) be an arbitrary object.
    Applying Universal Instantiation to \eqref{eq:statement-for-the-intersection-of-T1-substituted} yields,
    \begin{equation}
        \label{eq:instantiated-statement-for-the-intersection-of-T1-substituted}
        x_{0} \in \Intersect T_{1}
        \iff
        \forall y \, \mleft( \exists z \, \mleft( z \in S \land y = A \union z \mright) \implies x_{0} \in y \mright)
    \end{equation}
    Now we transform \eqref{eq:instantiated-statement-for-the-intersection-of-T1-substituted} into various logically
    equivalent forms:
    \begin{equation*}
        \begin{aligned}
            &\forall y \, \mleft( \exists z \, \mleft( z \in S \land y = A \union z \mright) \implies x_{0} \in y \mright)\\
            &\forall y \, \mleft( \lnot \, \exists z \, \mleft( z \in S \land y = A \union z \mright) \lor x_{0} \in y \mright)\\
            &\forall y \, \mleft( \forall z \, \mleft( z \notin S \lor y \neq A \union z \mright) \lor x_{0} \in y \mright)\\
            &\forall y \, \forall z \, \mleft( z \notin S \lor y \neq A \union z \lor x_{0} \in y \mright)\\
            &\forall z \, \forall y \, \mleft( z \notin S \lor y \neq A \union z \lor x_{0} \in y \mright)\\
            &\forall z \, \mleft( z \notin S \lor \forall y \, \mleft( y \neq A \union z \lor x_{0} \in y \mright) \mright)\\
            &\forall z \, \mleft( z \notin S \lor \forall y \, \mleft( y = A \union z \implies x_{0} \in y \mright) \mright)\\
            &\forall z \, \mleft( z \notin S \lor x_{0} \in A \union z \mright)\\
            &\forall z \, \mleft( z \notin S \lor \mleft( x_{0} \in A \lor x_{0} \in z \mright) \mright)\\
            &\forall z \, \mleft( \mleft( z \notin S \lor x_{0} \in A \mright) \lor x_{0} \in z \mright)\\
            &\forall z \, \mleft( \mleft( x_{0} \in A \lor z \notin S \mright) \lor x_{0} \in z \mright)\\
            &\forall z \, \mleft( x_{0} \in A \lor \mleft( z \notin S \lor x_{0} \in z \mright) \mright)\\
            &x_{0} \in A \lor \forall z \, \mleft( z \notin S \lor x_{0} \in  z \mright)\\
            &x_{0} \in A \lor \forall z \, \mleft( z \in S \implies x_{0} \in z \mright)\\
        \end{aligned}
        \text{  }
        \begin{aligned}
            &\quad \text{from \eqref{eq:instantiated-statement-for-the-intersection-of-T1-substituted}}\\
            &\quad \text{Definition of Conditional}\\
            &\quad \text{De Morgan Laws}\\
            &\quad \text{\cref{logic:grouping-disjuncts-into-quantified-statement-universal}}\\
            &\quad \text{Logical Equivalence}\\
            &\quad \text{\cref{logic:grouping-disjuncts-into-quantified-statement-universal}}\\
            &\quad \text{Definition of Conditional}\\
            &\quad \text{\cref{logic:using-universal-quantifier-to-express-the-fact-that-an-object-satisfies-a-predicate}}\\
            &\quad \text{\cref{thm:union-in-terms-of-logical-or}}\\
            &\quad \text{\( \lor \)-Associativity}\\
            &\quad \text{\( \lor \)-Commutativity}\\
            &\quad \text{\( \lor \)-Associativity}\\
            &\quad \text{\cref{logic:grouping-disjuncts-into-quantified-statement-universal}}\\
            &\quad \text{Definition of Conditional}\\
        \end{aligned}
    \end{equation*}
    Thus, we have,
    \begin{equation}
        \label{eq:final-statement-of-intersection-T1}
        x_{0} \in \Intersect T_{1}
        \iff
        x_{0} \in A \lor \forall z \, \mleft( z \in S \implies x_{0} \in z \mright)
    \end{equation}
    From \cref{thm:union-in-terms-of-logical-or,dfn:intersection-of-sets},
    \begin{align}
        \label{eq:final-statement-for-A-union-intersection-S}
        x_{0} \in A \union \Intersect S
        &\iff
        x_{0} \in A \lor x \in \Intersect S \notag \\
        &\iff
        x_{0} \in A \lor \forall z \, \mleft( z \in S \implies x_{0} \in z \mright)
    \end{align}
    Comparing \eqref{eq:final-statement-of-intersection-T1} and \eqref{eq:final-statement-for-A-union-intersection-S},
    it is clear that,
    \begin{equation*}
        x_{0} \in A \union \Intersect S \iff x_{0} \in \Intersect T_{1}
    \end{equation*}
    Since \( x_{0} \) is arbitrary, we conclude from Universal Generalization that,
    \begin{equation*}
        \forall x \, \mleft( x \in A \union \Intersect S \iff x \in \Intersect T_{1} \mright)
    \end{equation*}
    By the \nameref{axm:axiom-of-extensionality}, it follows that
    \begin{equation*}
        A \union \Intersect S = \Intersect T_{1}
    \end{equation*}
    The second equality can be similarly established. Again, we start on the right-hand side, this time using
    \cref{dfn:union-of-sets}.
    \begin{equation}
        \label{eq:statement-for-the-union-of-T2}
        \forall x \, \mleft( x \in \Union T_{2}
        \iff
        \exists y \, \mleft( y \in T_{2} \land x \in y \mright)\mright)
    \end{equation}
    Applying \eqref{eq:define-T2} to \eqref{eq:statement-for-the-union-of-T2},
    \begin{equation}
        \label{eq:statement-for-the-union-of-T2-substituted}
        \forall x \, \mleft( x \in \Union T_{2} \iff \exists y \, \mleft( \exists z \, \mleft( z \in S \land y = A \intersect z \mright) \land x \in y \mright) \mright)
    \end{equation}
    Let \( x_{0} \) be an arbitrary object. Instantiating \eqref{eq:statement-for-the-union-of-T2-substituted} yields,
    \begin{equation}
        \label{eq:instantiated-statement-for-the-union-of-T2-substituted}
        x_{0} \in \Union T_{2}
        \iff
        \exists y \, \mleft( \exists z \, \mleft( z \in S \land y = A \intersect z \mright) \land x_{0} \in y \mright)
    \end{equation}
    Now we begin transforming \eqref{eq:instantiated-statement-for-the-union-of-T2-substituted} into its equivalent forms:
    \begin{equation*}
        \begin{aligned}
            &\exists y \, \mleft( \exists z \, \mleft( z \in S \land y = A \intersect z \mright) \,{\land}\,\, x_{0} \in y \mright)\\
            &\exists y \, \exists z \, \mleft( z \in S \land y = A \intersect z \land x_{0} \in y \mright)\\
            &\exists z \, \exists y \, \mleft( z \in S \land y = A \intersect z \land x_{0} \in y \mright)\\
            &\exists z \, \mleft( z \in S \land \exists y \, \mleft( y = A \intersect z \land x_{0} \in y \mright) \mright)\\
            &\exists z \, \mleft( z \in S \land x_{0} \in A \intersect z \mright)\\
            &\exists z \, \mleft( z \in S \,\,{\land}\, \mleft( x_{0} \in A \land x_{0} \in z \mright)\mright)\\
            &\exists z \, \mleft( \mleft( z \in S \land x_{0} \in A \mright) \,{\land}\,\, x_{0} \in z \mright)\\
            &\exists z \, \mleft( \mleft( x_{0} \in A \land z \in S \mright) \,{\land}\,\, x_{0} \in z \mright)\\
            &\exists z \, \mleft( x_{0} \in A \,\,{\land}\, \mleft( z \in S \land x_{0} \in z \mright)\mright)\\
            &x_{0} \in A \land \exists z \, \mleft( z \in S \land x_{0} \in z \mright)\\
        \end{aligned}
        \text{  }
        \begin{aligned}
            &\quad \text{from \eqref{eq:instantiated-statement-for-the-union-of-T2-substituted}}\\
            &\quad \text{\cref{logic:grouping-conjuncts-into-quantified-statement-existential}}\\
            &\quad \text{Logical Equivalence}\\
            &\quad \text{\cref{logic:grouping-conjuncts-into-quantified-statement-existential}}\\
            &\quad \text{\cref{logic:using-existential-quantifier-to-express-the-fact-that-an-object-satisfies-a-predicate}}\\
            &\quad \text{\cref{thm:intersection-in-terms-of-logical-and}}\\
            &\quad \text{\( \land \)-Associativity}\\
            &\quad \text{\( \land \)-Commutativitiy}\\
            &\quad \text{\( \land \)-Associativity}\\
            &\quad \text{\cref{logic:using-existential-quantifier-to-express-the-fact-that-an-object-satisfies-a-predicate}}\\
        \end{aligned}
    \end{equation*}
    Thus, we have
    \begin{equation}
        \label{eq:final-statement-for-the-union-of-T2}
        x_{0} \in \Union T_{2}
        \iff
        x_{0} \in A \land \exists z \, \mleft( z \in S \land x_{0} \in z \mright)
    \end{equation}
    And from \cref{thm:intersection-in-terms-of-logical-and,dfn:union-of-sets},
    \begin{align}
        \label{eq:final-statement-of-A-intersect-union-S}
        x_{0} \in A \intersect \Union S
        &\iff
        x_{0} \in A \land x_{0} \in \Union S \notag \\
        &\iff
        x_{0} \in A \land \exists z \, \mleft( z \in S \land x_{0} \in z \mright)
    \end{align}
    Comparing \eqref{eq:final-statement-for-the-union-of-T2} and \eqref{eq:final-statement-of-A-intersect-union-S},
    it is clear that,
    \begin{equation*}
        x_{0} \in A \intersect \Union S \iff x_{0} \in \Union T_{2}
    \end{equation*}
    Since \( x_{0} \) is arbitrary, we conclude,
    \begin{equation*}
        \forall x \, \mleft( x \in A \intersect \Union S \iff x \in \Union T_{2}  \mright)
    \end{equation*}
    By the \nameref{axm:axiom-of-extensionality}, it follows that,
    \begin{equation*}
        A \intersect \Union S = \Union T_{2} \qedhere
    \end{equation*}
\end{proof}

\begin{thm}[Distributive Laws]
    \label{set-identities:distributive-laws}
    For any three sets \( A \), \( B \), and \( C \),
    \begin{gather*}
        A \union \mleft( B \intersect C \mright)
        = \mleft( A \union B \mright) \intersect \mleft( A \union C \mright)\\
        A \intersect \mleft( B \union C \mright)
        = \mleft( A \intersect B \mright) \union \mleft( A \intersect C \mright)
    \end{gather*}
\end{thm}

\begin{proof}
    This is a special case of \cref{set-identities:generalized-distributive-laws} with
    \( S = \mleft\{ B, C \mright\} \).
    Thus, by \cref{set-identities:generalized-distributive-laws,dfn:intersection-of-sets} we have,
    \begin{align*}
        A \union \mleft( B \intersect C \mright)
        &=
        A \union \Intersect \Big\{ B, C \Big\} \\
        &=
        \Intersect \setbuilderbig{x}{\exists z \, \mleft( z \in \mleft\{ B,C \mright\} \land x = A \union z \mright)}\\
        &=
        \Intersect \Big\{ A \union B, A \union C \Big\} \\
        &=
        \mleft( A \union B \mright) \intersect \mleft( A \union C \mright)
    \end{align*}
    Similarly,
    \begin{align*}
        A \intersect \mleft( B \union C \mright)
        &=
        A \intersect \Union \Big\{ B, C \Big\}\\
        &=
        \Union \setbuilderbig{x}{\exists z \, \mleft( z \in \mleft\{ B,C \mright\} \land x = A \intersect z \mright)}\\
        &=
        \Union \Big\{ A \intersect B, A \intersect C \Big\}\\
        &=
        \mleft( A \intersect B \mright) \union \mleft( A \intersect C \mright) \qedhere
    \end{align*}
\end{proof}

\begin{thm}[Absorption Laws of \( \union  \) and \( \intersect  \)]
    \label{thm:absorption-laws-of-union-and-intersection}
    For any two sets \( A \) and \( B \), the following equality holds:
    \begin{equation*}
        A \union \mleft( A \intersect B \mright) = A \intersect \mleft( A \union B \mright) = A
    \end{equation*}
\end{thm}

\begin{proof}
    Let \( A \) and \( B \) be arbitrary sets.
    It is trivial to prove that \( A \union \mleft( A \intersect B \mright) = A \intersect \mleft( A \union B \mright) \).
    By \cref{set-identities:distributive-laws}, we have,
    \begin{equation*}
        A \union \mleft( A \intersect B \mright) = \mleft( A \union A \mright) \intersect \mleft( A \union B \mright)
    \end{equation*}
    The \( \union  \) operation is idempotent (\cref{set-identities:idempotence-of-unions}), so \( A \union A = A \).
    \begin{equation*}
        \therefore \, A \union \mleft( A \intersect B \mright) = A \intersect \mleft( A \union B \mright)
    \end{equation*}
    Now, we simply have to prove \( A \intersect \mleft( A \union B \mright) = A \).
    \begin{align*}
        A \intersect \mleft( A \union B \mright)
        &=
        \mleft( A \union \emptyset \mright) \intersect \mleft( A \union B \mright) \\
        &=
        \big( A \union \mleft( B - B \mright) \big) \intersect \big( A \union B \big)\\
        &=
        A \union \big( \mleft( B - B \mright) \intersect B \big)\\
        &=
        A \union \mleft( \emptyset \intersect B \mright) \\
        &=
        A \union \emptyset \\
        &=
        A  \qedhere
    \end{align*}
\end{proof}

\begin{thm}[Right Distributivity of \( - \)]
    \label{set-identities:right-distributivity-of-difference}
    For any three sets \( A \), \( B \), and \( C \), the following
    equalities hold:
    \begin{gather*}
        \mleft( A \union B \mright) - C
        =
        \mleft( A - C \mright) \union \mleft( B - C \mright) \\
        \mleft( A \intersect B \mright) - C
        =
        \mleft( A - C \mright) \intersect \mleft( B - C \mright) \\
        \mleft( A - B \mright) - C
        =
        \mleft( A - C \mright) - \mleft( B - C \mright) \\
        \mleft( A \symdiff B \mright) - C
        =
        \mleft( A - C \mright) \symdiff \mleft( B - C \mright)
    \end{gather*}
\end{thm}

\begin{proof}
    The first two identities are simple to prove. It is evident from
    \cref{thm:union-in-terms-of-logical-or,dfn:difference} that for
    some arbitrary \( x_{0} \),
    \begin{align*}
        x_{0} \in \mleft( A \union B \mright) - C
        &\iff
        x_{0} \in \mleft( A \union B \mright) \land x_{0} \notin C \\
        &\iff
        \mleft( x_{0} \in A \lor x_{0} \in B \mright) \land x_{0} \notin C \\
        &\iff
        \mleft( x_{0} \in A \land x_{0} \notin C \mright) \lor
        \mleft( x_{0} \in B \land x_{0} \notin C \mright) \\
        &\iff
        x_{0} \in \mleft( A - C \mright) \lor x_{0} \in \mleft( B - C \mright) \\
        &\iff
        x_{0} \in \mleft( A - C \mright) \union \mleft( B - C \mright)
    \end{align*}
    Since \( x_{0} \) is arbitrary, the formula can be generalized,
    \begin{equation*}
        \forall x \, \Big( x \in \mleft( A \union B \mright) - C
        \iff
        x \in \mleft( A - C \mright) \union \mleft( B - C \mright)  \Big)
    \end{equation*}
    By the \nameref{axm:axiom-of-extensionality}, \( \mleft( A \union B \mright) - C = \mleft( A - C \mright) \union \mleft( B - C \mright) \).

    \vspace{0.5cm}
    \noindent
    To prove the second identitiy, we follow a similar procedure. Let \( x_{0} \) be an arbitrary object. Thus,
    \begin{align*}
        x_{0} \in \mleft( A \intersect B \mright) - C
        &\iff
        x_{0} \in \mleft( A \intersect B \mright) \land x_{0} \notin C\\
        &\iff
        \mleft( x_{0} \in A \land x_{0} \in B \mright) \land x_{0} \notin C\\
        &\iff
        \mleft( x_{0} \in A \land x_{0} \in B \mright)
        \land
        \mleft( x_{0} \notin C \land x_{0} \notin C \mright)\\
        &\iff
        \mleft( x_{0} \in A \land x_{0} \notin C \mright)
        \land
        \mleft( x_{0} \in B \land x_{0} \notin C \mright)\\
        &\iff
        x_{0} \in \mleft( A - C \mright) \land x_{0} \in \mleft( B - C \mright)\\
        &\iff
        x_{0} \in \mleft( A - C \mright) \intersect \mleft( B - C \mright)
    \end{align*}
    The proof for the third identity is slightly more involved.
    We shall employ a neat trick from propositional logic that
    allows us to derive an equivalent formula from any formula by
    introducing a false disjunct. That is, for any two formulas
    \( \varphi \) and \( \Psi \), the following equivalence holds:
    \begin{equation*}
        \varphi \equiv \varphi \lor \mleft( \Psi \land \lnot \, \Psi \mright)
    \end{equation*}
    Thus,
    \begin{align*}
        x_{0} \in \mleft( A - B \mright) - C
        &\iff
        \mleft( x_{0} \in A \land x_{0} \notin B \mright) \land x_{0} \notin C \\
        &\iff
        \big( \mleft( x_{0} \in A \land x_{0} \notin B \mright) \land x_{0} \notin C \big) \lor
        \big( x_{0} \in C \land x_{0} \notin C \big)\\
        &\iff
        \big( \mleft( x_{0} \in A \land x_{0} \notin B \mright) \lor x_{0} \in C \big)
        \land
        x_{0} \notin C \\
        &\iff
        \big( \mleft( x_{0} \in A \lor x_{0} \in C \mright) \land \mleft( x_{0} \notin B \lor x_{0} \in C \mright) \big)
        \land
        x_{0} \notin C\\
        &\iff
        \big( \mleft( x_{0} \in A \lor x_{0} \in C \mright) \land x_{0} \notin C \big)
        \land
        \big( x_{0} \notin B \lor x_{0} \in C \big)\\
        &\iff
        \big( x_{0} \in A \land x_{0} \notin C \big)
        \land
        \lnot \, \big( x_{0} \in B \land x_{0} \notin C \big)\\
        &\iff
        x_{0} \in \mleft( A - C \mright) \land x_{0} \notin \mleft( B - C \mright)\\
        &\iff
        x_{0} \in \big( \mleft( A - C \mright) - \mleft( B - C \mright) \big)
    \end{align*}
    And finally, we prove the fourth identity using the previous three identities.
    Applying \cref{set-identities:symmetric-difference-in-terms-of-union-and-intersection},
    we get,
    \begin{equation*}
        \mleft( A \symdiff B \mright) - C = \Big( \mleft( A \union B \mright) - \mleft( A \intersect B \mright) \Big) - C
    \end{equation*}
    Using the third identity, the expression becomes,
    \begin{equation*}
        \mleft( A \symdiff B \mright) - C
        =
        \Big( \mleft( A \union B \mright) - C - \mleft( A \intersect B \mright) - C \Big)
    \end{equation*}
    Then apply the first two identities,
    \begin{equation*}
        \mleft( A \symdiff B \mright) - C
        =
        \Big( \mleft( A - C \mright) \union \mleft( B - C \mright) - \mleft( A - C \mright) \intersect \mleft( B - C \mright) \Big)
    \end{equation*}
    Now apply \cref{set-identities:symmetric-difference-in-terms-of-union-and-intersection} again to obtain,
    \begin{equation*}
        \mleft( A \symdiff B \mright) - C = \mleft( A - C \mright) \symdiff \mleft( B - C\mright) \qedhere
    \end{equation*}
\end{proof}

\begin{thm}
    \label{set-identities:associativity-of-intersection-difference}
    For any three sets \( A \), \( B \), and \( C \):
    \begin{equation*}
        A \intersect \mleft( B - C \mright)
        = \mleft( A \intersect B \mright) - C
        = \mleft( A - C \mright) \intersect \mleft( B - C \mright)
    \end{equation*}
\end{thm}

\begin{proof}
    The equality follows from the associativity of the logical-\( \land \) operator.
    Thus, for any arbitrary \( x_{0} \),
    \begin{align*}
        x_{0} \in A \intersect \mleft( B - C \mright)
        &\iff
        x_{0} \in A \land x_{0} \in B - C\\
        &\iff
        x_{0} \in A \land \mleft( x_{0} \in B \land x_{0} \notin C \mright)\\
        &\iff
        \mleft( x_{0} \in A \land x_{0} \in B \mright) \land x_{0} \notin C\\
        &\iff
        \mleft( x_{0} \in A \intersect B \mright) \land \mleft( x_{0} \notin C \mright)\\
        &\iff
        x_{0} \in \mleft( A \intersect B \mright) - C
    \end{align*}
    To prove that \( A \intersect \mleft( B \intersect C \mright) = \mleft( A - C \mright) \intersect \mleft( B - C \mright) \),
    we simply apply \cref{set-identities:right-distributivity-of-difference} to
    \( \mleft( A \intersect B \mright) - C \).
\end{proof}

\begin{thm}
    \label{set-identities:union-difference}
    For any three sets \( A \), \( B \), and \( C \), the following equality holds:
    \begin{equation*}
        A \union \mleft( B - C \mright) = \mleft( A \union B \mright) - \mleft( C - A \mright)
    \end{equation*}
\end{thm}

\begin{proof}
    Let \( x_{0} \) be an arbitrary object. By \cref{thm:union-in-terms-of-logical-or,dfn:difference}, we have,
    \begin{equation*}
        x_{0} \in A \union \mleft( B - C \mright)
        \iff
        x_{0} \in A \lor \mleft( x_{0} \in B \land x_{0} \notin C \mright)
    \end{equation*}
    Since \( \lor \) distributes over \( \land \),
    \begin{equation*}
        x_{0} \in A \union \mleft( B - C \mright)
        \iff
        \mleft( x_{0} \in A \lor x_{0} \in B \mright)
        \land
        \mleft( x_{0} \in A \lor x_{0} \notin C \mright)
    \end{equation*}
    Applying De Morgan Laws,
    \begin{align*}
        x_{0} \in A \union \mleft( B - C \mright)
        &\iff
        \mleft( x_{0} \in A \lor x_{0} \in B \mright)
        \land \lnot \,
        \mleft( x_{0} \in C \land x_{0} \notin A \mright)\\
        &\iff
        x_{0} \in \mleft( A \union B \mright)
        \land
        x_{0} \notin \mleft( C - A \mright)\\
        &\iff
        x_{0} \in \Big( \mleft( A \union B \mright) - \mleft( C - A \mright) \Big) \qedhere
    \end{align*}
\end{proof}

\begin{thm}[Distributivity of \( \intersect  \) over \( \symdiff  \)]
    \label{set-identities:distributivity-of-intersection-over-symmetric-difference}
    For any three sets \( A \), \( B \), and \( C \), the following equalities hold:
    \begin{gather*}
        A \intersect (B \symdiff C)
        =
        \mleft( A \intersect B \mright) \symdiff \mleft( A \intersect C \mright)\\
        \mleft( A \symdiff B \mright) \intersect C
        =
        \mleft( A \intersect C \mright) \symdiff \mleft( B \intersect C \mright)
    \end{gather*}
\end{thm}

\begin{proof}
    We start by proving the first equality and then use it to prove the second equality.
    Let \( x_{0} \) be an arbitrary object. By definition,
    \begin{align*}
        x_{0} \in A \intersect \mleft( B \symdiff C \mright)
        &\iff
        x_{0} \in A \land x_{0} \in \mleft( B \symdiff C \mright)\\
        &\iff
        x_{0} \in A \, \land \mleft( x_{0} \in B \land x_{0} \notin C \lor x_{0} \notin B \land x_{0} \in C \mright)
    \end{align*}
    Simplifying using the distributivity of logical operators,
    \begin{equation}
        \label{eq:working-expression}
        \mleft( x_{0} \in A \land x_{0} \in B \land x_{0} \notin C \mright)
        \lor
        \mleft( x_{0} \in A \land x_{0} \notin B \land x_{0} \in C \mright)
    \end{equation}
    For the next step, we utilize the fact that for any formula \( \varphi \),
    \begin{equation*}
        \varphi \equiv \varphi \lor \bot
    \end{equation*}
    For instance, the disjunct on the left-hand side of \eqref{eq:working-expression} is equivalent
    to,
    \begin{equation*}
        \mleft( x_{0} \in A \land x_{0} \in B \land x_{0} \notin A \mright)
        \lor
        \mleft( x_{0} \in A \land x_{0} \in B \land x_{0} \notin C \mright)
    \end{equation*}
    We can simply the formula even further using the distributivity of logical operators.
    \begin{equation}
        \label{eq:left-disjunct-final-form}
        \mleft( x_{0} \in A \land x_{0} \in B \land x_{0} \notin C \mright)
        \iff
        \mleft( x_{0} \in A \land x_{0} \in B \mright)
        \land
        \mleft( x_{0} \notin A \lor x_{0} \notin C \mright)
    \end{equation}
    The same can be done for the disjunct on the right-hand side of \eqref{eq:working-expression} this time using
    \( x_{0} \in A \land x_{0} \in C \land x_{0} \notin A\) as the contradictory formula.
    \begin{equation}
        \label{eq:right-disjunct-final-form}
        \mleft( x_{0} \in A \land x_{0} \notin B \land x_{0} \in C \mright)
        \iff
        \mleft( x_{0} \in A \land x_{0} \in C \mright)
        \land
        \mleft( x_{0} \notin A \lor x_{0} \notin B \mright)
    \end{equation}
    Applying \eqref{eq:left-disjunct-final-form} and \eqref{eq:right-disjunct-final-form} into
    \eqref{eq:working-expression} yields,
    \begin{multline*}
        \Big( \mleft( x_{0} \in A \land x_{0} \in B \mright)
        \land
        \mleft( x_{0} \notin A \lor x_{0} \notin C \mright)\Big)
        \lor \\
        \Big( \mleft( x_{0} \in A \land x_{0} \in C \mright)
        \land
        \mleft( x_{0} \notin A \lor x_{0} \notin B \mright)\Big)
    \end{multline*}
    which is equivalent to,
    \begin{equation}
        \label{eq:final-equivalent-form}
        \Big( x_{0} \in \mleft( A \intersect B \mright) \land x_{0} \notin \mleft( A \intersect C \mright)  \Big)
        \lor
        \Big( x_{0} \in \mleft( A \intersect C \mright) \land x_{0} \notin \mleft( A \intersect B \mright) \Big)
    \end{equation}
    And by \cref{dfn:symmetric-difference}, \eqref{eq:final-equivalent-form} is equivalent to
    \( x_{0} \in  \mleft( A \intersect B \mright) \symdiff \mleft( A \intersect C \mright) \).
    \begin{equation*}
        \therefore \, x_{0} \in A \intersect \mleft( B \symdiff C \mright)
        \iff
        x_{0} \in \mleft( A \intersect B \mright) \symdiff \mleft( A \intersect C \mright)
    \end{equation*}
    Since \( x_{0} \) is arbitrary,
    \( A \intersect \mleft( B \symdiff C \mright) = \mleft( A \intersect B \mright) \symdiff \mleft( A \intersect C \mright)\)
    follows from Universal Generalization and the \nameref{axm:axiom-of-extensionality}.

    \vspace{0.5cm}
    \noindent
    With the first equality at our disposal, the proof for the second equality becomes trivial.
    Due to the commutativity of \( \intersect \), the problem is reduced to the first equality:
    \begin{align*}
        \mleft( A \symdiff B \mright) \intersect C
        &=
        C \intersect \mleft( A \symdiff B \mright)\\
        &=
        \mleft( C \intersect A \mright) \symdiff \mleft( C \intersect B \mright) \\
        &=
        \mleft( A \intersect C \mright) \symdiff \mleft( B \intersect C \mright) \qedhere
    \end{align*}
\end{proof}
